\chapter{Grassmanniany, algebra Clifforda i spinory}
\label{ch:grass-spin}
\index{grupa spin@Grupa Spin}

W Rozdziale~\ref{ch:algebra-tensorowa} skonstruowaliśmy iloczyn
zewnętrzny i wiązki tensorowe. Teraz zadamy trzy kolejne pytania:
jak opisać \emph{przestrzeń wszystkich podprzestrzeni} danego wymiaru,
jak zmienić algebrę zewnętrzną, aby uwzględniała iloczyn skalarny,
i dlaczego obrót o $2\pi$ może zmieniać znak spinora.
Odpowiedzi prowadzą kolejno do grassmannianów, algebr Clifforda
i grup Spin. Przykłady w wymiarach $2$ i $3$ pozwolą nam
sprawdzać każdą konwencję znaku na konkretnych macierzach.

\section{Grassmannian jako przestrzeń podprzestrzeni}
\label{sec:grass-atlas}

Przez $\mathbb K$ oznaczamy w tej sekcji $\R$ albo $\C$.
Jeśli $0\leq k\leq n$, \emph{grassmannian}
$\operatorname{Gr}_k(\mathbb K^n)$ to zbiór wszystkich
$k$-wymiarowych podprzestrzeni liniowych $\mathbb K^n$.
Jego \emph{punktem} jest cała podprzestrzeń $L$,
a nie wektor w $L$. Na przykład
$\operatorname{Gr}_1(\R^n)=\mathbb{RP}^{n-1}$
z Przykładu~\ref{prz:RPn}, a
$\operatorname{Gr}_1(\C^n)=\mathbb{CP}^{n-1}$.
Dla $k=0$ i $k=n$ grassmannian jest jednym punktem.
Dla dowolnej $n$-wymiarowej przestrzeni $V$ nad $\mathbb K$
piszemy analogicznie $\operatorname{Gr}_k(V)$; wybór bazy
identyfikuje go z powyższym modelem, lecz konstrukcje nie zależą od tej bazy.

\begin{twierdzenie}[Mapy wykresowe grassmannianu]
\label{thm:grass-atlas}
Niech $\mathbb K^n=P\oplus Q$, gdzie $\dim_{\mathbb K}P=k$.
Zbiór
\[
 U_P=\{L\in\operatorname{Gr}_k(\mathbb K^n):
           \pi_P|_L:L\longrightarrow P\text{ jest izomorfizmem}\}
\]
można jednoznacznie opisać odwzorowaniami liniowymi
$A:P\to Q$: punktowi $A$ odpowiada jego
\emph{wykres} $\Gamma(A)=\{u+Au:u\in P\}$.
Mapy $\Gamma(A)\mapsto A$ tworzą gładki atlas
grassmannianu. Jego wymiar nad $\mathbb K$ wynosi
$k(n-k)$. W przypadku zespolonym atlas jest holomorficzny,
a rzeczywisty wymiar wynosi $2k(n-k)$.
Oznaczenie $U_P$ zależy również od wybranego dopełnienia $Q$;
pomijamy je w indeksie tylko dla skrócenia zapisu.
\end{twierdzenie}
\begin{proof}
Jeśli $L=\Gamma(A)$, to $\pi_P(u+Au)=u$, więc
$\pi_P|_L$ jest izomorfizmem. W drugą stronę, dla
$L\in U_P$ określmy
\[
 A_L:=\pi_Q|_L\circ(\pi_P|_L)^{-1}:P\longrightarrow Q.
\]
Każdy $v\in L$ ma postać $u+\pi_Qv=u+A_Lu$
dla $u=\pi_Pv$. Stąd $L=\Gamma(A_L)$; dowiedliśmy
zarówno istnienia, jak i jednoznaczności parametru.

Sprawdźmy przejścia między mapami, a nie tylko ich
liczbę. Niech $\mathbb K^n=P'\oplus Q'$ będzie drugim
rozkładem. Dla $A:P\to Q$ połóżmy
$T_A:P\to\mathbb K^n$, $T_A(u)=u+Au$.
Wykres $\Gamma(A)$ należy do $U_{P'}$ dokładnie wtedy,
gdy $D_A:=\pi_{P'}T_A:P\to P'$ jest odwracalny.
Po wybraniu baz $P,P'$ wyznacznik macierzy $D_A$ zależy
wielomianowo od współczynników $A$;
warunek $\det D_A\ne0$ jest więc otwarty. W nowej mapie
parametr ma postać
\begin{equation}\label{eq:grass-przejscie}
 A'=\pi_{Q'}T_A\circ(\pi_{P'}T_A)^{-1}:P'\longrightarrow Q'.
\end{equation}
Odwracanie macierzy jest gładkie na zbiorze macierzy
odwracalnych, a nad $\C$ jest tam także holomorficzne
(wzór przez dopełnienia algebraiczne i wyznacznik).
Każde $L$ leży w $U_L$, bo dla rozkładu
$\mathbb K^n=L\oplus Q$ otrzymujemy $A_L=0$.
Atlasy pokrywają zatem cały zbiór. Współczynniki
$A\in\operatorname{Hom}_{\mathbb K}(P,Q)$ tworzą
przestrzeń $\mathbb K^{(n-k)k}$, co daje wymiar.

Definiujemy topologię przez atlas: zbiór jest otwarty, gdy
jego obraz w każdej z map wykresowych jest otwarty.
Otwartość dziedzin przejść i ich ciągłość zapewniają, że
same $U_P$ są otwarte, a mapy są homeomorfizmami.
Uzasadnijmy również pozostałe własności topologiczne.
Wybierzmy standardowy iloczyn skalarny, a nad $\C$ hermitowski;
gwiazdka oznacza operator sprzężony. W tym kroku
ograniczmy się do rozkładów ortogonalnych $Q=P^\perp$;
takie mapy wciąż pokrywają cały grassmannian
(dla danego $L$ bierzemy $P=L$).
Wzór~\eqref{eq:grass-przejscie} zapewnia zgodność
z pozostałymi mapami. Każdemu $L$
przypiszmy macierz rzutowania ortogonalnego $\Pi_L$.
Jest ona wyznaczona jednoznacznie przez $L$, a jej
współczynniki w mapie wykresowej są gładkie:
\[
 \Pi_{\Gamma(A)}
  =T_A(T_A^*T_A)^{-1}T_A^*.
\]
Odwrót na $U_P$ jest również ciągły w macierzach
rzutowania, ponieważ dla $L\in U_P$
\[
 A_L=\pi_Q\Pi_L|_P\,(\pi_P\Pi_L|_P)^{-1}.
\]
Ostatnia odwrotność istnieje, gdyż przy ortogonalnym
rzutowaniu $\pi_P$ odwzorowanie $\Pi_L|_P:P\to L$
jest sprzężeniem izomorfizmu
$\pi_P|_L:L\to P$. Nasza topologia zgadza się więc
z topologią podzbioru macierzy rzutowania: $U_P$ jest w tym
modelu zbiorem, gdzie $\pi_P\Pi_L|_P$ jest odwracalne,
a więc zbiorem otwartym. Oba podane wzory są ciągłe na tych mapach.
Zbiór $\Pi^*=\Pi$, $\Pi^2=\Pi$, $\operatorname{tr}\Pi=k$
jest domknięty i ograniczony w skończenie wymiarowej
przestrzeni macierzy, a więc zwarty. Istotnie, samosprzężone
idempotenty mają wartości własne $0,1$, więc norma operatorowa
jest co najwyżej $1$, a ślad jest rangą. Są zatem dokładnie
rzutami na $k$-płaszczyzny. Ten model jest także
Hausdorffa i ma przeliczalną bazę. Atlas definiuje
zatem rozmaitość w sensie Rozdziału~\ref{ch:rozmaitosci}.
\end{proof}

\begin{figure}[!htbp]
\centering
\begin{tikzpicture}[>=Stealth,line cap=round,line join=round,scale=1.02]
 \filldraw[fill=manifoldblue!11,draw=manifoldblue!70!black,thick]
 (-2.55,-.3)--(1.45,-.3)--(2.55,.3)--(-1.45,.3)--cycle;
 \filldraw[fill=manifoldgreen!14,draw=manifoldgreen!75!black,thick]
 (-2.55,-1.15)--(1.45,.25)--(2.55,1.15)--(-1.45,-.25)--cycle;
 \draw[->,manifoldgold!85!black,very thick] (0,0)--(0,1.75)
  node[above,font=\small] {$Q$};
 \fill[black!75] (0,0) circle (2pt) node[below left,font=\small] {$0$};
 \fill[manifoldblue] (1.675,.15) circle (2pt)
  node[below right,fill=white,inner sep=1pt,font=\small] {$u\in P$};
 \fill[manifoldgreen!80!black] (1.675,.715) circle (2pt)
  node[above right,fill=white,inner sep=1pt,font=\small] {$u+Au\in\Gamma(A)$};
 \draw[->,manifoldred,thick] (1.675,.15)--(1.675,.715)
  node[midway,left,font=\small] {$Au$};
 \node[manifoldblue!70!black,font=\small\bfseries] at (-1.2,.52) {$P$};
 \node[manifoldgreen!65!black,font=\small\bfseries] at (-2.5,-1.48) {$\Gamma(A)$};
\end{tikzpicture}
\caption{Mapa wykresowa grassmannianu. Przestrzeń $P$ jest
punktem odniesienia; pobliska płaszczyzna $\Gamma(A)$ powstaje,
gdy nad $u\in P$ odkładamy przemieszczenie $Au\in Q$.}
\label{fig:grass-wykres}
\end{figure}

\begin{przyklad}[Mała mapa, a nie cały grassmannian]
Na $\operatorname{Gr}_1(\R^2)=\mathbb{RP}^1$ ustalmy
$P=\R e_1$ i $Q=\R e_2$. Prosta
$L_a=\operatorname{span}(e_1+ae_2)$ ma współrzędną
$a\in\R$. Brakuje dokładnie jednej prostej, $\R e_2$:
rzutowanie na $P$ jest na niej zerowe. Zatem mapa
wykresowa jest lokalna, podobnie jak mapa stereograficzna
nie pokrywa całej sfery.
\end{przyklad}

\section{Przestrzeń styczna i wiązka tautologiczna}
\label{sec:grass-styczna}

Wektor styczny do grassmannianu opisuje \emph{pierwszy
rząd zmiany podprzestrzeni}. Wektory należące już do
danej podprzestrzeni nie zmieniają jej położenia:
istotna jest część wychodząca poza nią.

\begin{twierdzenie}[Styczne jako odwzorowania do ilorazu]
\label{thm:grass-styczna}
Dla $L\in\operatorname{Gr}_k(V)$ istnieje kanoniczny
izomorfizm
\begin{equation}\label{eq:grass-styczna}
 T_L\operatorname{Gr}_k(V)
  \cong\operatorname{Hom}_{\mathbb K}(L,V/L).
\end{equation}
Nad $\C$ prawa strona jest zespoloną przestrzenią
styczną; jej podłożem rzeczywistym jest rzeczywista
przestrzeń styczna.
\end{twierdzenie}
\begin{proof}
Niech $L(t)$ będzie gładką krzywą podprzestrzeni,
$L(0)=L$. Dla $v\in L$ wybierzmy gładką krzywą
$v(t)\in L(t)$ z $v(0)=v$. Definiujemy
$D_{L'(0)}(v):=[v'(0)]\in V/L$, gdzie nawias oznacza
klasę modulo $L$. Pokażmy niezależność od wyboru
$v(t)$. Różnica dwóch takich krzywych $w(t)\in L(t)$
ma $w(0)=0$. W gładkiej lokalnej bazie $s_i(t)$
zapisuje się $w(t)=\sum_i a_i(t)s_i(t)$ z
$a_i(0)=0$, zatem
$w'(0)=\sum_i a_i'(0)s_i(0)\in L$. Klasa
$[v'(0)]$ nie zmienia się. Zmiana parametryzacji
krzywej nie jest jedyną możliwą zmianą przedstawiciela wektora
stycznego; niezależność od dowolnej krzywej o tym samym wektorze
stycznym sprawdzimy teraz w mapie.

Ustalmy teraz uzupełnienie $V=L\oplus Q$.
W mapie z Twierdzenia~\ref{thm:grass-atlas} mamy
$L(t)=\Gamma(A(t))$ i $A(0)=0$.
Wybieramy $v(t)=v+A(t)v$. Wtedy
\[
 D_{L'(0)}(v)=[A'(0)v]\in V/L.
\]
Tak samo $s_i(t)=e_i+A(t)e_i$ dla bazy $(e_i)$ przestrzeni $L$
daje gładką bazę używaną powyżej. Wzór zależy tylko od $A'(0)$,
czyli współrzędnych wektora stycznego, a nie od dalszego przebiegu krzywej.
Jeśli $D_{L'(0)}=0$, to $A'(0)v\in Q\cap L=0$ dla każdego $v$,
więc $A'(0)=0$ i wektor styczny jest zerowy. To dowodzi iniektywności.
Utożsamienie $Q\cong V/L$ pokazuje, że
$D_{L'(0)}$ jest liniowe w $v$. Każdy
$D\in\operatorname{Hom}(L,V/L)$ można otrzymać:
wybierzmy odpowiadające mu odwzorowanie $A:L\to Q$
i krzywą $L(t)=\Gamma(tA)$. Jej pochodna daje
właśnie $D$. Konstrukcja rozpoczęła się od
krzywej i klas modulo $L$, więc otrzymany
izomorfizm nie zależy od wybranego $Q$.
\end{proof}

\begin{definicja}[Wiązka tautologiczna]
\index{wiazka tautologiczna@Wiązka tautologiczna}
\label{def:grass-taut}
Nad $\operatorname{Gr}_k(V)$ określamy
\emph{wiązkę tautologiczną}
\[
 \gamma^k
 :=\{(L,v)\in\operatorname{Gr}_k(V)\times V:v\in L\}
 \longrightarrow\operatorname{Gr}_k(V),\qquad (L,v)\mapsto L.
\]
Jej włókno w punkcie $L$ jest dosłownie podprzestrzenią
$L\subseteq V$.
\end{definicja}

\begin{stwierdzenie}[Lokalność i wiązka ilorazowa]
\label{prop:grass-wiazki}
Wiązka $\gamma^k$ ma rząd $k$. Jest podwiązką
wiązki trywialnej $\operatorname{Gr}_k(V)\times V$.
Istnieje też wiązka ilorazowa $\mathcal Q$ z
włóknem $\mathcal Q_L=V/L$, rzędu $n-k$,
oraz kanoniczny dokładny ciąg wiązek
\begin{equation}\label{eq:grass-ciag}
 0\longrightarrow\gamma^k
 \longrightarrow\operatorname{Gr}_k(V)\times V
 \longrightarrow\mathcal Q\longrightarrow0.
\end{equation}
Po wybraniu iloczynu skalarnego otrzymujemy
$\mathcal Q\cong(\gamma^k)^\perp$.
\end{stwierdzenie}
\begin{proof}
Na $U_P$ każdy element $L=\Gamma(A)$ zawiera
dokładnie jeden wektor $u+Au$ dla danego $u\in P$.
Odwzorowanie $(L,u)\mapsto(L,u+Au)$ jest zatem lokalną
trywializacją $U_P\times P\cong\gamma^k|_{U_P}$;
jej odwrotność to $(L,v)\mapsto(L,\pi_Pv)$.
Oba odwzorowania są gładkie, a na włóknach liniowe.
W tym samym rozkładzie $V=P\oplus Q$ klasa
$[w]\in V/\Gamma(A)$ ma jednoznacznego
reprezentanta w $Q$: od $w$ odejmujemy
$u+Au\in\Gamma(A)$, gdzie $u=\pi_Pw$.
Stąd lokalnie $\mathcal Q|_{U_P}\cong U_P\times Q$.
Na każdym włóknie włączenie $L\hookrightarrow V$
jest iniekcją, a naturalne odwzorowanie $V\to V/L$ jest
surjekcją z jądrem $L$, co wyjaśnia \emph{dokładność}
ciągu~\eqref{eq:grass-ciag}. Iloczyn skalarny
identyfikuje klasę $[v]$ z $(\id-\Pi_L)v\in L^\perp$.
Rzut zależy gładko od $L$, więc jest to izomorfizm wiązek,
a nie tylko rodzina izomorfizmów włókien.
\end{proof}

\begin{uwaga}[Jeszcze jedno zastosowanie Rozdziału 3]
Twierdzenie~\ref{thm:grass-styczna} i
izomorfizm tensorowy~\ref{thm:tensor-hom} dają
\[
 T\operatorname{Gr}_k(V)\cong
       (\gamma^k)^*\otimes\mathcal Q.
\]
Włókno po prawej stronie w $L$ to
$L^*\otimes(V/L)\cong\operatorname{Hom}(L,V/L)$.
Kolejność czynników jest tu ważna: źródłem odwzorowania
jest $L$, a jego wartości należą do $V/L$.
Zgodność tych izomorfizmów z gładkością również widać w mapie:
przy $L=\Gamma(A)$ przyrost $\dot A$ wysyła $u+Au$ na
klasę $\dot A u$. W trywializacji ilorazu
$[w]\mapsto\pi_Qw-A\pi_Pw$ ma on po prostu macierz $\dot A$.
Zatem izomorfizmy włókien sklejają się w gładki izomorfizm wiązek.
\end{uwaga}

\begin{przyklad}[Proste w płaszczyźnie i wstęga Möbiusa]
Na $\operatorname{Gr}_1(\R^2)\cong\mathbb{RP}^1\cong S^1$
wiązka $\gamma^1$ jest wstęgą Möbiusa z
Przykładu~\ref{prz:mobius}. Istotnie,
$L(t)=\R(\cos t,\sin t)$, $0\leq t\leq\pi$,
opisuje pełne okrążenie bazy: $L(0)=L(\pi)$.
Wybrany lokalnie wektor jednostkowy
$(\cos t,\sin t)\in L(t)$ wraca jako jego
przeciwieństwo. To właśnie sklejanie włókna
ze zmianą znaku.
\end{przyklad}

\section{Orientacja, współrzędne Plückera i odwzorowanie Gaussa}
\label{sec:grass-plucker}

\begin{definicja}[Grassmannian zorientowany]
\index{grassmannian zorientowany@Grassmannian zorientowany}
Punkt $\operatorname{Gr}^+_k(\R^n)$ to para $(L,o)$:
$k$-płaszczyzna $L$ wraz z wyborem jej orientacji $o$.
Zapomnienie orientacji daje dwukrotne nakrycie
$\operatorname{Gr}^+_k(\R^n)\to\operatorname{Gr}_k(\R^n)$
dla $0<k<n$: nad każdą płaszczyzną są dokładnie
dwie możliwe orientacje. W mapie wykresowej
orientację przenosimy z $P$ na $\Gamma(A)$ przez
$u\mapsto u+Au$, więc lokalnie nakrycie jest
sumą dwóch kopii tej samej mapy.
\end{definicja}

\begin{przyklad}[Płaszczyzny w $\R^3$]
Płaszczyźnie $L\in\operatorname{Gr}_2(\R^3)$
odpowiada jej nieukierunkowana prosta normalna
$L^\perp$, a zatem
$\operatorname{Gr}_2(\R^3)\cong\mathbb{RP}^2$.
Gdy $L$ ma orientację, wybierzmy dodatnią bazę
ortonormalną $(u,v)$ i połóżmy $n=u\times v$.
Inna dodatnia baza różni się obrotem w $L$,
więc daje ten sam $n$. W drugą stronę $n\in S^2$
wyznacza $L=n^\perp$ i orientację, dla której
$(u,v,n)$ jest dodatnią bazą $\R^3$. Otrzymujemy
\[
 \operatorname{Gr}^+_2(\R^3)\cong S^2,\qquad
 \operatorname{Gr}_2(\R^3)\cong S^2/(n\sim -n).
\]
Zmiana orientacji płaszczyzny zamienia $n$ z $-n$.
Dla ustalenia gładkości zauważmy, że normalna z lokalnej dodatniej
ramki ortonormalnej jest gładka. W drugą stronę rzut
na $n^\perp$ ma macierz $I-nn^T$, gładką w $n$;
lokalny wybór orientacji jest stały na każdym z dwóch płatów.
Są to więc dyfeomorfizmy, nie tylko bijekcje zbiorów.
\end{przyklad}

\begin{figure}[!htbp]
\centering
\begin{tikzpicture}[>=Stealth,line cap=round,scale=1.02]
 \filldraw[fill=manifoldblue!10,draw=manifoldblue!65!black,thick]
 (-4.6,-.52)--(-1.3,-.52)--(-.45,.08)--(-3.75,.08)--cycle;
 \draw[->,very thick,manifoldred] (-2.55,-.22)--(-2.55,1.05)
   node[above,font=\small] {$n$};
 \draw[->,thick,manifoldgold!85!black,dashed]
   (-2.55,-.22)--(-2.55,-1.49)
   node[below,font=\small] {$-n$};
 \node[font=\small,manifoldblue!80!black] at (-1.2,.35) {$L$};
 \draw[->,manifoldblue!70!black] (-3.35,-.31)--(-2.45,-.31)
   node[midway,below,font=\scriptsize] {$u$};
 \draw[->,manifoldblue!70!black] (-3.35,-.31)--(-2.85,.01)
   node[near end,above,font=\scriptsize] {$v$};
 \shade[inner color=white,outer color=manifoldblue!43] (2.1,0) circle (1.12);
 \draw[manifoldblue!75!black,thick] (2.1,0) circle (1.12);
 \draw[manifoldgreen!65,thin] (.98,0) arc[start angle=180,end angle=360,x radius=1.12,y radius=.27];
 \draw[manifoldgreen!65,thin,dashed] (3.22,0)
   arc[start angle=0,end angle=180,x radius=1.12,y radius=.27];
 \fill[manifoldred] (2.1,1.12) circle (2.3pt)
   node[above,font=\small] {$n$};
 \fill[manifoldgold!85!black] (2.1,-1.12) circle (2.3pt)
   node[below,font=\small] {$-n$};
 \draw[->,manifoldblue!75!black,thick] (-.05,.15)--(.68,.15)
   node[midway,above,font=\small] {$\perp$};
 \node[font=\small] at (2.1,-1.65) {$S^2\longrightarrow\mathbb{RP}^2$};
\end{tikzpicture}
\caption{Orientacja płaszczyzny w $\R^3$ wybiera jedną z
dwóch jednostkowych normalnych. Zapomnienie orientacji
utożsamia punkty antypodyczne $n$ i $-n$ na sferze.}
\label{fig:grass-normalne}
\end{figure}

\begin{stwierdzenie}[Bazy ortonormalne i przestrzeń ilorazowa]
\label{prop:grass-stiefel}
Niech $V_k(\R^n)$ będzie zbiorem uporządkowanych
ortonormalnych $k$-ramek $(v_1,\ldots,v_k)$.
Odwzorowanie $(v_i)\mapsto\operatorname{span}(v_i)$ ma
włókna izomorficzne z $\mathrm O(k)$.
Ponadto, jako gładkie przestrzenie,
\[
 \operatorname{Gr}_k(\R^n)
   \cong\mathrm O(n)/(\mathrm O(k)\times\mathrm O(n-k)),
\quad
 \operatorname{Gr}_k(\C^n)
   \cong\mathrm U(n)/(\mathrm U(k)\times\mathrm U(n-k)).
\]
\end{stwierdzenie}
\begin{proof}
Dwie ramki wyznaczają tę samą $k$-płaszczyznę wtedy
i tylko wtedy, gdy zmiana bazy między nimi należy do
$\mathrm O(k)$. Wykonując ortonormalizację Grama--Schmidta
na wektorach $u_i+Au_i$ z mapy wykresowej, dostajemy
gładki lokalny wybór ramki; włókna są więc lokalnie
kopiami $\mathrm O(k)$. Grupa $\mathrm O(n)$ działa
tranzytywnie na płaszczyznach: wybieramy bazy ortonormalne
dwóch płaszczyzn, uzupełniamy je do baz $\R^n$ i bierzemy
odwzorowanie przenoszące pierwszą pełną bazę na drugą.
Element zachowujący standardowe $\R^k$ zachowuje też
jego dopełnienie prostopadłe, czyli ma postać blokową
$\mathrm{diag}(B,C)$ z $B\in\mathrm O(k)$,
$C\in\mathrm O(n-k)$. To daje pierwszy iloraz.
Wyjaśnijmy gładką strukturę ilorazu. W rozkładzie ortogonalnym
$P\oplus Q$ uzupełniamy ramkę wykresu przez ramkę
$\Gamma(-A^*)\subset Q\oplus P$ jego dopełnienia prostopadłego:
jej wektory mają postać $-A^*w+w$, $w\in Q$.
Ortonormalizacja obu ramek daje macierz $s(A)\in\mathrm O(n)$
zależną gładko od $A$. Każda macierz ortogonalna nad tą mapą
ma jednoznaczną postać $s(A)\operatorname{diag}(B,C)$.
Parametr $A$ odzyskujemy z obrazu pierwszych $k$ kolumn,
a bloki $B,C$ z iloczynu $s(A)^{-1}O$. Daje to lokalny
produkt i jego gładką odwrotność; po podzieleniu przez bloki
pozostaje dokładnie mapa wykresowa grassmannianu.
Argument dla $\C$ jest taki sam, z bazami unitarnymi,
sprzężeniem hermitowskim i grupą $\mathrm U(n)$.
\end{proof}

Iloczyn zewnętrzny z Sekcji~\ref{sec:algebra-zewnetrzna}
pozwala zapisać $k$-płaszczyznę jednym wielowektorem,
określonym z dokładnością do niezerowego skalara.

\begin{twierdzenie}[Zanurzenie Plückera]
\label{thm:grass-plucker}
Jeśli $(v_1,\ldots,v_k)$ jest bazą $L$, to wzór
\begin{equation}\label{eq:plucker}
 \mathrm{Pl}(L):=[v_1\wedge\cdots\wedge v_k]
 \in\mathbb P(\Lambda^k V)
\end{equation}
określa gładkie zanurzenie
$\operatorname{Gr}_k(V)\hookrightarrow\mathbb P(\Lambda^k V)$.
Nawias $[\xi]$ oznacza prostą $\mathbb K\xi$,
$\xi\ne0$. Obraz składa się dokładnie z klas
niezerowych \emph{wielowektorów prostych}.
\end{twierdzenie}
\begin{proof}
Zmiana bazy $(v_i)$ macierzą $C$ mnoży iloczyn
zewnętrzny przez $\det C\ne0$ na mocy
Twierdzenia~\ref{thm:det-wlasnosci}, więc klasa
projektowa jest dobrze określona. Jeżeli
$\xi=v_1\wedge\cdots\wedge v_k\ne0$, to
\[
 L=\{w\in V:w\wedge\xi=0\}.
\]
Wektory z $L$ dają zero przez powtórzenie czynnika.
Jeśli $w\notin L$, rodzina $(v_1,\ldots,v_k,w)$
jest liniowo niezależna; po uzupełnieniu do bazy
Twierdzenie~\ref{thm:wedge-baza} mówi, że jej
iloczyn zewnętrzny jest niezerowy. Zatem
$\mathrm{Pl}(L)$ wyznacza $L$ jednoznacznie.

Pokażmy, skąd bierze się gładki odwrotny wzór na obrazie.
W mapie $V=P\oplus Q$ wybierzmy bazy
$e_1,\ldots,e_k$ przestrzeni $P$ oraz
$e_{k+1},\ldots,e_n$ przestrzeni $Q$ i zapiszmy
$Ae_j=\sum_{a=1}^{n-k}a_{aj}e_{k+a}$.
W wielowektorze
$\xi_A=(e_1+Ae_1)\wedge\cdots\wedge(e_k+Ae_k)$
współczynnik $p_{1\cdots k}$ wynosi $1$.
Współczynnik przy
$e_1\wedge\cdots\wedge\widehat{e_j}\wedge\cdots
\wedge e_k\wedge e_{k+a}$ wynosi
$(-1)^{k-j}a_{aj}$: przesuwamy $e_{k+a}$
przez $k-j$ późniejszych czynników.
Z ilorazów tych współczynników przez $p_{1\cdots k}$
odzyskujemy wszystkie $a_{aj}$.
W ogólnej mapie projektowej dzielimy przez
odpowiedni niezerowy współczynnik $p_I$.
Odwzorowanie~\eqref{eq:plucker} jest więc gładkie.
Wzory odzyskujące $a_{aj}$ są określone na całym otwartym
zbiorze przestrzeni projektowej $p_{1\cdots k}\ne0$, nie tylko
na obrazie. Ich złożenie z $\mathrm{Pl}$ jest mapą identyczną
w parametrach $A$; różniczkowanie daje lewą odwrotność różniczki
$\mathrm{Pl}$, która jest zatem injektywna.
Otrzymaliśmy immersję. Ponieważ grassmannian jest zwarty,
a przestrzeń projektowa Hausdorffa, ciągła injekcja jest
homeomorfizmem na obraz. To daje zanurzenie i domkniętość obrazu.
\end{proof}

\begin{przyklad}[Relacja Plückera dla $\operatorname{Gr}_2(\R^4)$]
Zapiszmy $\xi=\sum_{1\leq i<j\leq4}p_{ij}e_i\wedge e_j$.
Rozwinięcie $\xi\wedge\xi$ daje
\begin{equation}\label{eq:plucker-24}
 \xi\wedge\xi=
 2(p_{12}p_{34}-p_{13}p_{24}+p_{14}p_{23})
   e_1\wedge e_2\wedge e_3\wedge e_4.
\end{equation}
Jeżeli $\xi=u\wedge v$, to $\xi\wedge\xi=0$,
więc współczynniki spełniają relację w nawiasie.
Odwrotnie, dla $\xi\ne0$ z taką relacją możemy
przestawić współrzędne tak, by $p_{12}\ne0$.
Po podzieleniu przez $p_{12}$ ustawiamy w wykresie
$e_1+ae_3+ce_4$, $e_2+be_3+de_4$ liczby
$b=p_{13}/p_{12}$, $d=p_{14}/p_{12}$,
$a=-p_{23}/p_{12}$, $c=-p_{24}/p_{12}$.
Pięć współczynników się zgadza z definicji,
a szósty jest równy $ad-bc=p_{34}/p_{12}$
właśnie dzięki relacji. To dowodzi również
wystarczalności. Na przykład
$e_1\wedge e_2+e_3\wedge e_4$ ma
$p_{12}p_{34}=1$ i pozostałe wyrazy zerowe,
więc nie jest wielowektorem prostym, zgodnie
z wcześniejszym Przykładem w
Sekcji~\ref{sec:baza-zewnetrzna}.
\end{przyklad}

\begin{twierdzenie}[Odwzorowanie Gaussa i wiązka styczna]
\label{thm:grass-gauss}
Niech $M^m\subset\R^N$ będzie gładką podrozmaitością.
Przypisanie
$G(p)=T_pM\in\operatorname{Gr}_m(\R^N)$
jest gładkim \emph{odwzorowaniem Gaussa}. Wiązki
$TM$ i normalna $\nu M$ są odpowiednio
cofnięciami $G^*\gamma^m$ i $G^*\mathcal Q$
z~\eqref{eq:grass-ciag}.
\end{twierdzenie}
\begin{proof}
W mapie podrozmaitości zapisujemy jej parametryzację
$F:U\subset\R^m\to\R^N$. Wektory
$\partial_iF(x)$ tworzą gładką bazę $T_{F(x)}M$.
W pobliżu ustalonego punktu pewien minor $m\times m$
macierzy $\dd F$ pozostaje niezerowy. W odpowiadającym
mu rozkładzie $\R^N=P\oplus Q$ przestrzeń
$T_{F(x)}M$ jest wykresem macierzy
$A(x)=D_QF(x)(D_PF(x))^{-1}$,
gładkiej w $x$. Stąd $G$ jest gładka.
Włókno $(G^*\gamma^m)_p$ to z definicji
$\gamma^m_{G(p)}=T_pM$; lokalne ramki po obu
stronach są tymi samymi $\partial_iF$.
Włókno $(G^*\mathcal Q)_p$ to
$\R^N/T_pM$, czyli $\nu_pM$ z
Przykładu o wiązce normalnej w Sekcji o podrozmaitościach
Rozdziału~\ref{ch:rozmaitosci}.
Zgodność lokalnych trywializacji daje oba
izomorfizmy wiązek.
\end{proof}

\begin{uwaga}[Dlaczego wiązka tautologiczna jest uniwersalna?]
Na zwartej rozmaitości każdą gładką wiązkę $E$
rzędu $k$ można umieścić w trywialnej wiązce
$M\times\R^N$ dla pewnego $N$.
Oto konstrukcja. Wybieramy skończone pokrycie
trywializujące $U_i$, jego mniejsze podpokrycie
oraz funkcje gładkie $\beta_i$ o nośnikach w $U_i$,
takie że dla każdego $p$ co najmniej jedna
$\beta_i(p)$ jest niezerowa. Jeśli $[v]_i\in\R^k$
oznacza współrzędne $v\in E_p$ na $U_i$,
definiujemy różnowartościowe odwzorowanie liniowe w $v$
\[
 j_p(v)=\bigl(\beta_1(p)[v]_1,\ldots,
                  \beta_r(p)[v]_r\bigr)\in\R^{kr},
\]
przyjmując wartość $0$ dla bloku poza $U_i$.
Podparcie $\beta_i$ wewnątrz $U_i$ zapewnia
gładkość po przedłużeniu zerem; co najmniej
jeden niezerowy blok zapewnia iniektywność $j_p$.
Obrazy $j_p(E_p)$ zmieniają się gładko, ponieważ
lokalnie tworzą kolumny macierzy gładkiej i
stałego rzędu $k$, a więc definiują odwzorowanie
$f:M\to\operatorname{Gr}_k(\R^{kr})$.
Wprost z definicji włókien mamy
$E\cong f^*\gamma^k$. Ta konkretna konstrukcja
wyjaśnia użyteczność grassmannianów dla topologii
wiązek; nie potrzebuje jeszcze teorii klas
charakterystycznych.
\end{uwaga}

\section{Algebra Clifforda: iloczyn i dowód wymiaru}
\label{sec:clifford}

Od tego miejsca $V$ jest rzeczywistą przestrzenią
euklidesową wymiaru $n$ z iloczynem skalarnym
$\langle\cdot,\cdot\rangle$ i normą
$\|v\|^2=\langle v,v\rangle$. Algebra zewnętrzna
pamięta zorientowane pola: $v\wedge v=0$.
Algebra Clifforda pamięta ponadto \emph{długość}:
kwadrat wektora jest skalarem $-\|v\|^2$.
Przyjmujemy konsekwentnie ten znak minus;
w literaturze spotyka się również konwencję przeciwną.

\begin{definicja}[Algebra Clifforda]
\index{algebra clifforda@Algebra Clifforda}
\label{def:clifford}
Niech $I_q$ będzie ideałem obustronnym algebry
tensorowej $T(V)$ z Definicji~\ref{def:algebra-tensorowa},
generowanym przez
$v\otimes v+\|v\|^2\,1$ dla wszystkich $v\in V$.
\emph{Algebra Clifforda} to
\[
 \operatorname{Cl}(V):=T(V)/I_q.
\]
Obraz $v\in V$ oznaczamy nadal $v$,
a mnożenie w ilorazie zapisujemy bez symbolu $\otimes$.
Z definicji $v^2=-\|v\|^2\,1$.
\end{definicja}

Po usunięciu z relacji składnika $\|v\|^2\,1$
otrzymujemy algebrę zewnętrzną
(Definicja~\ref{def:algebra-zewnetrzna}).
Tak więc obie algebry startują z tej samej algebry
tensorowej; zmieniamy tylko relację, przez którą
dzielimy. W algebrze Clifforda nie można zachować
zwykłego stopnia tensorowego: relacja utożsamia
wyraz stopnia $2$ ze skalarem stopnia $0$.
Pozostaje rozróżnienie stopnia \emph{parzystego}
i \emph{nieparzystego}, bo oba wyrazy relacji
mają tę samą parzystość.

\begin{twierdzenie}[Relacja Clifforda i własność uniwersalna]
\label{thm:cliff-relacje}
Dla $v,w\in V$ w $\operatorname{Cl}(V)$ zachodzi
\begin{equation}\label{eq:cliff-relacja}
 vw+wv=-2\langle v,w\rangle\,1.
\end{equation}
Jeżeli $\mathcal A$ jest łączną algebrą rzeczywistą
z jedynką, a $f:V\to\mathcal A$ jest liniowa i spełnia
$f(v)^2=-\|v\|^2 1_{\mathcal A}$ dla każdego $v$,
to istnieje dokładnie jeden homomorfizm algebr
$\widehat f:\operatorname{Cl}(V)\to\mathcal A$
przedłużający $f$.
\end{twierdzenie}
\begin{proof}
Rozwijając $(v+w)^2$ na dwa sposoby, otrzymujemy
\[
 v^2+vw+wv+w^2=-\|v+w\|^2\,1
  =-(\|v\|^2+\|w\|^2+2\langle v,w\rangle)\,1.
\]
Po odjęciu $v^2=-\|v\|^2 1$ i
$w^2=-\|w\|^2 1$ zostaje~\eqref{eq:cliff-relacja}.
Z Twierdzenia~\ref{thm:algebra-tensorowa-universal}
odwzorowanie $f$ przedłuża się jednoznacznie do homomorfizmu
$F:T(V)\to\mathcal A$. Dla każdego generatora ideału
$I_q$ mamy
$F(v\otimes v+\|v\|^2 1)
 =f(v)^2+\|v\|^2 1_{\mathcal A}=0$.
Jądro homomorfizmu jest ideałem, więc zawiera $I_q$.
Wobec tego wzór $\widehat f([t])=F(t)$
nie zależy od przedstawiciela $t\in T(V)$
i wyznacza żądany homomorfizm na ilorazie.
Jednoznaczność wynika stąd, że $V$ i jedynka
generują $\operatorname{Cl}(V)$.
\end{proof}

Przy bazie ortonormalnej $(e_1,\ldots,e_n)$
relacja~\eqref{eq:cliff-relacja} mówi
\begin{equation}\label{eq:cliff-orto}
 e_i^2=-1,\qquad e_ie_j=-e_je_i\quad(i\ne j).
\end{equation}
Stąd każdy iloczyn generatorów można uporządkować
rosnąco, redukując powtórzone indeksy.
Trudniejszą częścią jest wykazanie, że tak otrzymane
wyrazy nie mają dodatkowych relacji.

\begin{twierdzenie}[Baza algebry Clifforda]
\label{thm:cliff-baza}
Dla bazy ortonormalnej $e_1,\ldots,e_n$ elementy
\[
 1,\qquad e_{i_1}\cdots e_{i_r}
 \quad(1\leq i_1<\cdots<i_r\leq n,\ 1\leq r\leq n)
\]
stanowią bazę $\operatorname{Cl}(V)$. Zatem
$\dim_\R\operatorname{Cl}(V)=2^n$.
Dla $n\geq1$ części parzysta i nieparzysta mają
po $2^{n-1}$ wymiarów.
\end{twierdzenie}
\begin{proof}
\emph{Rozpinanie.} Obrazy słów $v_1\cdots v_r$
rozpinają iloraz $T(V)/I_q$. Po rozwinięciu każdego
$v_j$ w bazie wystarczą słowa z liter $e_i$.
Równości~\eqref{eq:cliff-orto} pozwalają zamienić
każde dwie sąsiednie różne litery (ze zmianą znaku)
i usunąć parę równych liter ($e_i^2=-1$).
Powtarzając te operacje, dostajemy kombinację
iloczynów o rosnących, różnych indeksach.

\emph{Niezależność.} Skorzystamy z algebry zewnętrznej
$\Lambda(V)$ z Rozdziału~\ref{ch:algebra-tensorowa}.
Na niej określmy dwie operacje:
\[
 \varepsilon(v)\omega:=v\wedge\omega,\qquad
 \iota_v(w_1\wedge\cdots\wedge w_r)
 :=\sum_{j=1}^r(-1)^{j-1}
   \langle v,w_j\rangle\,
   w_1\wedge\cdots\wedge\widehat{w_j}
       \wedge\cdots\wedge w_r .
\]
Kreska nad $w_j$ oznacza pominięcie tego czynnika;
przyjmujemy $\iota_v(1)=0$.
Wzór na $\iota_v$ jest wieloliniowy i alternujący
w $w_j$: przy zamianie dwóch sąsiednich argumentów składniki
pomijające inny argument zmieniają znak przez antysymetrię iloczynu
zewnętrznego, a dwa pozostałe zamieniają się miejscami z przeciwnymi
znakami. Zatem własność uniwersalna
z Twierdzenia~\ref{thm:wedge-multilinear}
zapewnia, że operator jest dobrze określony.
Bezpośrednio z tego wzoru otrzymujemy
\[
 \iota_v(w\wedge\omega)
  =\langle v,w\rangle\omega-w\wedge\iota_v\omega.
\]
Pierwszy składnik powstaje przy usunięciu pierwszego
czynnika; pozostałe mają o jeden minus więcej niż
w $\iota_v\omega$. Usunięcie dwóch czynników
w dwóch przeciwnych kolejnościach daje składniki
o przeciwnych znakach, zatem $\iota_v^2=0$.
Również $\varepsilon(v)^2=0$, bo $v\wedge v=0$.
Jeżeli $c(v):=\varepsilon(v)-\iota_v$, to
\[
 c(v)^2
 =-\bigl(\varepsilon(v)\iota_v
          +\iota_v\varepsilon(v)\bigr)
 =-\|v\|^2\id_{\Lambda(V)}.
\]
Własność uniwersalna z poprzedniego twierdzenia
daje więc reprezentację
$c:\operatorname{Cl}(V)\to\operatorname{End}_\R(\Lambda V)$.
Przyłóżmy ją do $1\in\Lambda^0V$.
Dla różnych ortonormalnych $e_{i_j}$ wszystkie
kontrakcje pojawiające się podczas kolejnego
działania znikają: nowy $e_{i_j}$ jest prostopadły
do wcześniej utworzonych czynników. Stąd
\[
 c(e_{i_1}\cdots e_{i_r})\,1
    =e_{i_1}\wedge\cdots\wedge e_{i_r},\qquad
 c(1)\,1=1.
\]
Prawa strona przebiega bazę $\Lambda(V)$ z
Twierdzenia~\ref{thm:wedge-baza}. Jeśli kombinacja
uporządkowanych słów znikałaby w algebrze Clifforda,
jej działanie na $1$ dawałoby zerową kombinację
\emph{różnych} wektorów tej bazy. Wszystkie
współczynniki muszą więc być zerowe.
Słów o $r$ indeksach jest $\binom nr$;
sumując po $r$ otrzymujemy $2^n$.
Wreszcie bijekcja $I\mapsto I\triangle\{1\}$
zamienia podzbiory parzystej i nieparzystej
liczności, co dowodzi równości ich liczby.
\end{proof}

\begin{uwaga}[Podobne przestrzenie, różne mnożenia]
\label{rem:cliff-zewnetrzna}
Istnieje kanoniczny izomorfizm \emph{przestrzeni liniowych}
$\Lambda V\cong\operatorname{Cl}(V)$: wysyła
$v_1\wedge\cdots\wedge v_r$ na antysymetryzację
$\frac1{r!}\sum_{\sigma\in S_r}\operatorname{sgn}(\sigma)
 v_{\sigma(1)}\cdots v_{\sigma(r)}$.
Jest dobrze określony, bo wyrażenie jest alternujące;
na bazowych wektorach ortonormalnych daje
$e_{i_1}\cdots e_{i_r}$, więc poprzednie twierdzenie
zapewnia odwracalność. Nie jest to jednak izomorfizm
algebr: w $\Lambda V$ mamy $v\wedge v=0$,
a w $\operatorname{Cl}(V)$ mamy
$v^2=-\|v\|^2$. To wyjaśnia, dlaczego obie
przestrzenie mają wymiar $2^n$, choć kodują
inne mnożenia.
\end{uwaga}

\begin{przyklad}[Wymiary $0$, $1$ i $2$]
Dla $V=0$ otrzymujemy $\operatorname{Cl}(0)=\R$.
W wymiarze $1$, przy $e_1^2=-1$, każdy element
ma postać $a+be_1$ i mnoży się jak liczba
$a+bi\in\C$. Stąd $\operatorname{Cl}(\R)\cong\C$
\emph{jako algebra rzeczywista}.
W wymiarze $2$ baza $(1,e_1,e_2,e_1e_2)$
spełnia dokładnie relacje kwaternionów:
$e_1\leftrightarrow i$, $e_2\leftrightarrow j$,
$e_1e_2\leftrightarrow k$, zatem
$\operatorname{Cl}(\R^2)\cong\mathbb H$.
Na przykład $(e_1+2e_2)^2=-5$, podczas gdy
$(e_1+2e_2)\wedge(e_1+2e_2)=0$
w algebrze zewnętrznej. Porównujemy więc kwadraty tego samego wektora.
\end{przyklad}

\begin{stwierdzenie}[Iloczyn wektorów: skalar i dwuwektor]
\label{prop:cliff-wedge}
Przez powyższy izomorfizm przestrzeni liniowych
iloczyn Clifforda dwóch wektorów rozkłada się na
\begin{equation}\label{eq:cliff-rozdzial}
 vw=-\langle v,w\rangle
       +\tfrac12(vw-wv)
 \quad\longleftrightarrow\quad
 -\langle v,w\rangle+v\wedge w.
\end{equation}
\end{stwierdzenie}
\begin{proof}
Suma $vw+wv$ jest równa
$-2\langle v,w\rangle$ z~\eqref{eq:cliff-relacja}.
Dodając połowę tej sumy do połowy różnicy,
odzyskujemy $vw$. Antysymetryzacja
$v\wedge w\mapsto\tfrac12(vw-wv)$
jest przypadkiem stopnia $2$ opisanym w
poprzedniej uwadze.
\end{proof}

\begin{figure}[!htbp]
\centering
\begin{tikzpicture}[>=Stealth,line cap=round,scale=1.0]
 \coordinate (O) at (0,0);
 \coordinate (U) at (3.0,0);
 \coordinate (W) at ({3*cos(47)},{3*sin(47)});
 \coordinate (UW) at ({3+3*cos(47)},{3*sin(47)});
 \fill[manifoldgreen!16] (O)--(U)--(UW)--(W)--cycle;
 \draw[manifoldgreen!60!black,thick] (U)--(UW)--(W);
 \draw[->,manifoldblue,very thick] (O)--(U)
  node[pos=.86,below,font=\small] {$u=e_1$};
 \draw[->,manifoldred,very thick] (O)--(W)
  node[pos=.78,above left,font=\small] {$v$};
 \draw[densely dashed,manifoldred!55] (W)--({3*cos(47)},0);
 \draw[->,manifoldred!70!black] (0,-.3)--({3*cos(47)},-.3)
  node[midway,below,font=\small] {$\cos\alpha$};
 \draw[manifoldgold!80!black,thick] (.58,0)
  arc[start angle=0,end angle=47,radius=.58];
 \node[manifoldgold!80!black,font=\small] at (.83,.32) {$\alpha$};
 \node[manifoldgreen!70!black,font=\small] at (3.3,1.3)
  {$u\wedge v=\sin\alpha\,e_1\wedge e_2$};
\end{tikzpicture}
\caption{Dla jednostkowych $u=e_1$ i
$v=\cos\alpha\,e_1+\sin\alpha\,e_2$ iloczyn Clifforda
$uv=-\cos\alpha+\sin\alpha\,e_1e_2$ łączy
ujemny iloczyn skalarny ze zorientowanym polem.}
\label{fig:cliff-iloczyn}
\end{figure}

\section{Od odbić do grupy Spin}
\label{sec:spin-grupa}

Ustalmy najpierw dwie operacje na algebrze Clifforda.
Zmiana każdego $v$ na $-v$ zachowuje relacje definiujące,
więc daje automorfizm $\alpha$, z $\alpha^2=\id$.
Jego przestrzenie własne dla $+1$ i $-1$ są rozpięte
odpowiednio przez słowa parzyste i nieparzyste; dowód bazy
pokazuje, że ich suma jest prosta.
Drugą operacją jest \emph{rewersja}, oznaczana tyldą:
\[
 \widetilde{v_1\cdots v_r}=v_r\cdots v_1,\qquad
 \widetilde 1=1,
\]
przedłużona liniowo. Jest dobrze określona na ilorazie:
odwracanie słów w $T(V)$ ustala każdy generator
$v\otimes v+\|v\|^2 1$, a element $a(v\otimes v+\|v\|^2 1)b$
wysyła na $\widetilde b(v\otimes v+\|v\|^2 1)\widetilde a$,
nadal należący do ideału. Zatem
$\widetilde{st}=\widetilde t\,\widetilde s$ i
$\widetilde{\widetilde s}=s$, niezależnie od zapisu $s$ jako sumy słów.

Odbicie względem hiperpłaszczyzny prostopadłej
do jednostkowego $u\in V$ ma wzór
\begin{equation}\label{eq:spin-odbicie}
 r_u(v)=v-2\langle u,v\rangle u.
\end{equation}
Rozkład $v=\langle u,v\rangle u+w$, $w\perp u$,
pokazuje, że $r_u$ odwraca składową równoległą do $u$
i zostawia składową $w$. Jego wyznacznik jest równy $-1$.
Relacja Clifforda daje zaskakujący opis tego odbicia:
ponieważ $u^2=-1$, mamy $u^{-1}=-u$ i
\[
 -uvu^{-1}=uvu
 =-\langle u,v\rangle u+w
 =r_u(v).
\]

\begin{definicja}[Grupa Pin]
\label{def:pin-grupa}
\index{grupa pin@Grupa Pin}
Dla euklidesowej przestrzeni $V$ grupą $\operatorname{Pin}(V)$
nazywamy podgrupę odwracalnych elementów $\operatorname{Cl}(V)$
złożoną ze wszystkich skończonych iloczynów wektorów
jednostkowych; dopuszczamy pusty iloczyn $1$.
W odróżnieniu od grupy Spin dopuszczamy także nieparzystą
liczbę czynników. Dla $s\in\operatorname{Pin}(V)$ określamy
\[
 \rho_{\mathrm{Pin}}(s)(v):=\alpha(s)v s^{-1},
 \qquad v\in V,
\]
gdzie $\alpha$ zmienia znak każdego wektora, jak wyżej.
\end{definicja}

Wzór jest skręconym sprzężeniem: dla jednostkowego $u$
daje $\rho_{\mathrm{Pin}}(u)(v)=-uvu^{-1}=r_u(v)$.
Zatem pojedynczy wektor reprezentuje odbicie, a jego parzysty
iloczyn reprezentuje obrót. Na przykład w $\R^2$ element
$e_1\in\operatorname{Pin}(2)$ daje odbicie w prostej
$e_1^\perp$, podczas gdy $e_1e_2$ daje obrót o pół obiegu.
Przy przyjętej relacji $u^2=-\|u\|^2$ grupa ta bywa
oznaczana $\operatorname{Pin}^{-}(n)$; znak rozróżnia dwie
konwencje algebry Clifforda.

\begin{definicja}[Grupa Spin]
\index{grupa spin@Grupa Spin}
\label{def:spin-grupa}
Niech $n=\dim V\geq2$. \emph{Grupą spinową}
$\operatorname{Spin}(V)$ nazywamy podgrupę
odwracalnych elementów $\operatorname{Cl}(V)$,
składającą się z iloczynów \emph{parzystej}
liczby wektorów jednostkowych:
\[
 \operatorname{Spin}(V)
 :=\{u_1\cdots u_{2r}:\|u_j\|=1,\ r\geq0\}.
\]
Pusty iloczyn to $1$. Skrót
$\operatorname{Spin}(n)$ oznacza
$\operatorname{Spin}(\R^n)$ ze standardową metryką.
Dla $s\in\operatorname{Spin}(V)$ określamy
$\rho(s)v:=svs^{-1}$.
Topologię na grupie bierzemy z jej włączenia do skończenie
wymiarowej przestrzeni $\operatorname{Cl}(V)$. Gładką strukturę
zgodną z tym włączeniem uzasadnimy w dowodzie nakrycia.
\end{definicja}

Zbiór w definicji jest grupą: iloczyn dwóch słów
parzystych jest parzysty, a odwrotność słowa
$u_1\cdots u_{2r}$ to
$(-u_{2r})\cdots(-u_1)$, ponownie słowo parzyste
złożone z wektorów jednostkowych.
Każdy jego czynnik ma dobrze określoną odwrotność.

\begin{twierdzenie}[Dwukrotne nakrycie grupy obrotów]
\label{thm:spin-nakrycie}
Przypisanie $\rho$ jest surjektywnym homomorfizmem
\[
 \rho:\operatorname{Spin}(n)\longrightarrow\mathrm{SO}(n),
 \qquad\ker\rho=\{1,-1\}.
\]
Jest gładkim dwukrotnym nakryciem: każdy obrót
ma dokładnie dwa podniesienia $s$ i $-s$.
\end{twierdzenie}
\begin{proof}
\emph{Dlaczego obrazem jest obrót?}
Dla słowa $s=u_1\cdots u_{2r}$ zastosujmy
wzór na odbicie do kolejnych czynników.
W odwzorowaniu $v\mapsto-u v u^{-1}$ pojawia się
jedno odbicie $r_u$; przy parzystej liczbie
czynników wszystkie minusy się znoszą i
\[
 \rho(s)=r_{u_1}\circ\cdots\circ r_{u_{2r}}.
\]
To złożenie zachowuje iloczyn skalarny i ma
wyznacznik $(-1)^{2r}=1$. Ponieważ
$\rho(st)(v)=s(tvt^{-1})s^{-1}$,
$\rho$ jest homomorfizmem do $\mathrm{SO}(n)$.

\emph{Dlaczego jest surjektywna?}
Pokażmy elementarnie, że każdą izometrię liniową
można zapisać jako iloczyn odbić. Weźmy bazę
ortonormalną $e_1,\ldots,e_n$ i $A\in\mathrm O(n)$.
Jeśli $Ae_1\ne e_1$, to dla
$u=(Ae_1-e_1)/\|Ae_1-e_1\|$ rachunek z
\eqref{eq:spin-odbicie} daje $r_u(Ae_1)=e_1$:
licznik iloczynu
$\langle Ae_1-e_1,Ae_1\rangle$ to
$1-\langle Ae_1,e_1\rangle$, a kwadrat
mianownika to $2(1-\langle Ae_1,e_1\rangle)$.
Jeśli $Ae_1=e_1$, pomijamy ten krok.
Nowa izometria zachowuje $e_1$, więc działa
na $e_1^\perp$. Powtarzamy argument dla $e_2$
w tym dopełnieniu, potem dla $e_3$, itd.
Po skończenie wielu krokach otrzymujemy
identyczność. Zapisujemy więc $A$ jako
iloczyn odbić. Ponieważ $\det A=1$, liczba
użytych odbić musi być parzysta; odpowiadający
im iloczyn jednostkowych $u_j$ jest szukanym
elementem $\operatorname{Spin}(n)$.

\emph{Dlaczego jądro ma tylko dwa elementy?}
Niech $s$ działa trywialnie na $V$. Wtedy
$se_j=e_js$ dla każdego wektora bazy ortonormalnej.
Rozwińmy $s$ w bazie z
Twierdzenia~\ref{thm:cliff-baza}, używając tylko
słów parzystej długości:
$s=\sum_{|I|\ {\rm parzyste}}a_I e_I$.
Dla parzystego niepustego $I$ i $j\in I$
relacja~\eqref{eq:cliff-orto} daje
$e_je_Ie_j^{-1}=-e_I$; gdy $j\notin I$,
otrzymujemy $+e_I$. Równość
$e_jse_j^{-1}=s$ dla \emph{wszystkich} $j$
wymusza więc $a_I=0$ dla każdego niepustego $I$.
Zatem $s=a\,1$. Aby ustalić $a$, odwróćmy
porządek wektorów w słowie:
$\widetilde s=u_{2r}\cdots u_1$.
Wtedy $s\widetilde s=(-1)^{2r}=1$,
bo wektory są jednostkowe i $u_i^2=-1$.
Dla skalarnego $s=a$ również $\widetilde s=a$,
stąd $a^2=1$. Oba znaki występują,
gdyż $-1=e_1e_1\in\operatorname{Spin}(n)$.

\emph{Dlaczego rzeczywiście jest to nakrycie?}
Pokażmy, że w pobliżu identyczności obroty mają
gładko wybrane podniesienie. Dla bliskich
jednostkowych $a,b$ połóżmy
\[
 R(a,b):=-\frac{a+b}{\|a+b\|}\,a
       =\frac{1-ba}{\|a+b\|}\in\operatorname{Spin}(n).
\]
Jest to iloczyn dwóch jednostkowych wektorów,
zależy gładko od $(a,b)$, spełnia $R(a,a)=1$
i przenosi $a$ na $b$: pierwsze odbicie w
hiperpłaszczyźnie do $a$ wysyła $a$ na $-a$,
drugie, o normalnej $(a+b)/\|a+b\|$, wysyła
$-a$ na $b$. Jeśli $A$ jest bliski identyczności,
stosujemy taki rotor, aby skorygować $Ae_1$ do
$e_1$. Otrzymany obrót zachowuje $e_1$.
W $e_1^\perp$ korygujemy następnie obraz
$e_2$, potem $e_3$, kończąc na $e_{n-1}$;
ostatni wektor jest ustalony dzięki wyznacznikowi
$1$. Każdy rotor zależy gładko od $A$
w niewielkim otoczeniu identyczności.
Odwrotność ich iloczynu daje gładką lokalną
sekcję $\sigma$ odwzorowania $\rho$.
Przy innym obrocie przenosimy to otoczenie
przez ustalony element z jego przeciwobrazu.
Jądro $\{\pm1\}$ oznacza, że przeciwobraz
takiego dostatecznie małego otoczenia jest
sumą rozłącznych kopii obrazu $\sigma$ i $-\sigma$.
Są one otwarte w $\operatorname{Spin}(n)$:
blisko $1$ dwie wartości $\pm\sigma(A)$
odróżnia znak współczynnika skalarnego
w algebrze Clifforda, a blisko pozostałych
punktów korzystamy z mnożenia przez
ustalone podniesienie. To jest dokładnie definicja
dwukrotnego nakrycia.
W ten sposób otrzymujemy także atlas gładki: na każdym płacie
używamy współrzędnych z $\mathrm{SO}(n)$ przeniesionych przez $\rho$.
Jest on zgodny z włączeniem do algebry Clifforda. Istotnie,
$\sigma$ jest gładka jako mapa do tej algebry, a pochodna
$\rho\circ\sigma=\id$ jest identycznością. Do różniczkowania
$\rho$ w otoczeniu płata wystarczy gładki wzór
$s\mapsto(v\mapsto\operatorname{pr}_V(svs^{-1}))$
na otwartym zbiorze elementów odwracalnych; na grupie Spin
rzut $\operatorname{pr}_V$ niczego nie zmienia.
Wynika stąd injektywność różniczki $\sigma$, więc płaty są
lokalnie podrozmaitościami algebry Clifforda. Mnożenie jest
dwuliniowe, a odwracanie gładkie: można je obliczać przez
odwracanie macierzy lewego mnożenia. Obie operacje ograniczają
się więc do gładkich operacji grupy Spin, a $\rho$ jest lokalnym dyfeomorfizmem.
\end{proof}

\begin{wniosek}[Nakrycie grupy ortogonalnej przez Pin]
\label{cor:pin-nakrycie}
Dla $n\geq2$ odwzorowanie
\[
 \rho_{\mathrm{Pin}}:\operatorname{Pin}(n)\longrightarrow
 \mathrm O(n)
\]
jest surjektywnym homomorfizmem o jądrze $\{1,-1\}$.
Jest gładkim dwukrotnym nakryciem, a przeciwobraz
$\mathrm{SO}(n)$ jest równy $\operatorname{Spin}(n)$.
\end{wniosek}
\begin{proof}
Automorfizm $\alpha$ jest multiplikatywny, więc
$\rho_{\mathrm{Pin}}(st)=\rho_{\mathrm{Pin}}(s)
\rho_{\mathrm{Pin}}(t)$. Każdy czynnik jednostkowy
daje odbicie. Rozkład izometrii na odbicia wykonano
w dowodzie Twierdzenia~\ref{thm:spin-nakrycie}; zapewnia
on surjektywność. Wyznacznik obrazu iloczynu $r$
wektorów wynosi $(-1)^r$. Parzyste słowa tworzą
zatem dokładnie przeciwobraz $\mathrm{SO}(n)$, czyli
$\operatorname{Spin}(n)$. Element jądra musi być parzysty,
a jądro na $\operatorname{Spin}(n)$ obliczono w tym samym
twierdzeniu jako $\{\pm1\}$. Lokalnie nad składową
$\mathrm{SO}(n)$ mamy już dwa gładkie płaty nakrycia Spin.
Mnożenie przez ustalony wektor jednostkowy przenosi je
na drugą składową $\mathrm O(n)$, co dowodzi ostatniej tezy.
\end{proof}

\begin{przyklad}[Dlaczego pojawia się połowa kąta?]
\label{prz:spin2}
W płaszczyźnie $B=e_1e_2$ spełnia $B^2=-1$.
Ponieważ część parzysta ma bazę $(1,B)$,
jej elementy mają postać $a+bB$. Dla słów
należących do $\operatorname{Spin}(2)$ równość
$s\widetilde s=1$ z poprzedniego dowodu daje
$a^2+b^2=1$. Odwrotnie, jeśli
$a=\cos t$, $b=\sin t$, to
\[
 a+bB=(-\cos t\,e_1-\sin t\,e_2)e_1,
\]
czyli iloczyn dwóch jednostkowych wektorów.
Grupa $\operatorname{Spin}(2)$ jest zatem okręgiem.
Połóżmy
\[
 s_\theta=\cos(\theta/2)+\sin(\theta/2)B.
\]
Korzystając z $Be_1=e_2$ i $Be_2=-e_1$,
obliczamy bez pomijania kroku
\[
 s_\theta e_1s_\theta^{-1}
 =(\cos^2\tfrac\theta2-\sin^2\tfrac\theta2)e_1
   +2\sin\tfrac\theta2\cos\tfrac\theta2\,e_2
 =\cos\theta\,e_1+\sin\theta\,e_2.
\]
Analogicznie obraz $e_2$ to
$-\sin\theta\,e_1+\cos\theta\,e_2$.
Jeden pełny obrót ma zatem podniesienie
$s_{2\pi}=-1$, a dopiero $s_{4\pi}=1$.
Po utożsamieniu obu grup z okręgiem odwzorowanie
$\rho$ ma postać $z\mapsto z^2$.
\end{przyklad}

\begin{figure}[!htbp]
\centering
\begin{tikzpicture}[>=Stealth,line cap=round]
 \draw[manifoldblue!75!black,thick] (-2.45,0) circle (1.03);
 \draw[manifoldgreen!75!black,thick] (2.45,0) circle (1.03);
 \draw[manifoldred,very thick,postaction={decorate},
  decoration={markings,mark=at position .55 with {\arrow{Stealth}}}]
  (-1.42,0) arc[start angle=0,end angle=180,radius=1.03];
 \draw[manifoldgreen!75!black,very thick,postaction={decorate},
  decoration={markings,mark=at position .78 with {\arrow{Stealth}}}]
  (3.48,0) arc[start angle=0,end angle=360,radius=1.03];
 \fill[manifoldblue] (-1.42,0) circle (2.2pt)
  node[right,font=\small] {$1$};
 \fill[manifoldred] (-3.48,0) circle (2.2pt)
  node[left,font=\small] {$-1$};
 \fill[manifoldgreen!70!black] (3.48,0) circle (2.2pt)
  node[right,font=\small] {$\id$};
 \draw[->,thick,manifoldgold!85!black] (-.84,0)--(.84,0)
  node[midway,above,font=\small] {$\rho:z\mapsto z^2$};
 \node[font=\small,manifoldblue!75!black] at (-2.45,-1.52)
  {$\operatorname{Spin}(2)$: pół obiegu};
 \node[font=\small,manifoldgreen!70!black] at (2.45,-1.52)
  {$\mathrm{SO}(2)$: cały obrót};
\end{tikzpicture}
\caption{Droga od $1$ do $-1$ w $\operatorname{Spin}(2)$
rzutuje się na pełen obrót w $\mathrm{SO}(2)$.
Końce drogi są różne w grupie Spin, choć odpowiadają
temu samemu obrotowi.}
\label{fig:spin-dwukrotne}
\end{figure}

\begin{przyklad}[Trzy wymiary: kwaterniony]
\label{prz:spin3}
W $\operatorname{Cl}(\R^3)$ część parzysta ma bazę
$1,I=e_2e_3,J=e_3e_1,K=e_1e_2$.
Rachunek przy użyciu~\eqref{eq:cliff-orto} daje
$I^2=J^2=K^2=-1$ i $IJ=K$, $JK=I$, $KI=J$.
Jest to algebra kwaternionów $\mathbb H$.
Odwracanie kolejności czynników wysyła
$a+bI+cJ+dK$ na $a-bI-cJ-dK$,
więc $s\widetilde s=a^2+b^2+c^2+d^2$.
Każdy element $\operatorname{Spin}(3)$
jest zatem jednostkowym kwaternionem.
I odwrotnie, dowolny jednostkowy kwaternion
można zapisać jako $\cos t+\sin t\,B$,
gdzie $B$ jest jednostkowym dwuwektorem
w $\R^3$ (przypadki $\pm1$ są natychmiastowe).
Każdy taki $B$ jest iloczynem dwóch
prostopadłych wektorów jednostkowych. Aby to sprawdzić,
zapiszmy $B=b e_2e_3+c e_3e_1+d e_1e_2$ z $b^2+c^2+d^2=1$.
Wybierzmy jednostkowy $u\perp n=(b,c,d)$ i połóżmy
$v=n\times u$. Wtedy $u,v$ są ortonormalne i $u\times v=n$.
Współczynniki $u\wedge v$ w bazie
$(e_2\wedge e_3,e_3\wedge e_1,e_1\wedge e_2)$ są właśnie $(b,c,d)$,
więc $uv=B$. Z kolei
$\cos t+\sin t\,B$ jest także iloczynem dwóch
wektorów jednostkowych: jeśli $B=uv$ dla
$u\perp v$, to wprost z $vu=-uv$ i $u^2=-1$
otrzymujemy
\[
 (-\cos t\,u-\sin t\,v)u
     =\cos t+\sin t\,uv.
\]
Zatem $\operatorname{Spin}(3)\cong S^3$
jako grupa jednostkowych kwaternionów,
a $\rho:S^3\to\mathrm{SO}(3)$ utożsamia
parę $s,-s$.
\end{przyklad}

\section{Spinory: moduły algebry Clifforda}
\label{sec:spinory}

Wektor w $\R^n$ obraca się przez $\rho(s)\in\mathrm{SO}(n)$.
Spinor ma podlegać działaniu samego $s\in\operatorname{Spin}(n)$,
tak aby dwa podniesienia tego samego obrotu mogły
działać odmiennie. Do opisania wszystkich wymiarów
użyjemy liczb zespolonych. \emph{Zespolenie}
$\operatorname{Cl}^{\C}_n:=
\operatorname{Cl}(\R^n)\otimes_\R\C$
ma tę samą bazę uporządkowanych słów co algebra
rzeczywista, ale współczynniki zespolone.
Jego wymiar nad $\C$ wynosi $2^n$.

\begin{definicja}[Spinor zespolony]
\index{spinor zespolony@Spinor zespolony}
\label{def:spinor}
\emph{Moduł} algebry $A$ nad $\mathbb K$ to przestrzeń liniowa
$\Delta$ nad $\mathbb K$ z działaniem $A\times\Delta\to\Delta$
dwuliniowym nad $\mathbb K$ i zgodnym
z mnożeniem: $(ab)\psi=a(b\psi)$
oraz $1\psi=\psi$. Jest \emph{nieredukowalny},
jeśli $\Delta\ne0$ i nie ma właściwej niezerowej podprzestrzeni
zachowywanej przez wszystkie działania $a\in A$.
W tym rozdziale \emph{przestrzenią spinorów}
w wymiarze $n$ nazwiemy nieredukowalny zespolony
moduł $\Delta_n$ algebry $\operatorname{Cl}^{\C}_n$.
Jej element $\psi$ to spinor. Ograniczenie
działania do $\operatorname{Spin}(n)\subset
\operatorname{Cl}^{\C}_n$ jest
\emph{reprezentacją spinową}.
\end{definicja}

Słowo \emph{moduł} brzmi abstrakcyjnie, lecz tutaj
oznacza po prostu konkretne macierze $c(v)$
działające na kolumny $\psi$, z warunkiem
$c(v)^2=-\|v\|^2 I$. Działanie na całej algebrze
jest wtedy wyznaczone przez jej własność uniwersalną.
Przykładem modułu jest $\Lambda V$ z dowodu
Twierdzenia~\ref{thm:cliff-baza}, ale zwykle jest
on \emph{większy} niż nieredukowalna przestrzeń
spinorów. Nie należy więc utożsamiać automatycznie
wielowektora ze spinorem.

\begin{przyklad}[Macierze Pauliego]
\label{prz:pauli}
Niech
\[
 \sigma_1=\begin{pmatrix}0&1\\1&0\end{pmatrix},
 \quad
 \sigma_2=\begin{pmatrix}0&-i\\i&0\end{pmatrix},
 \quad
 \sigma_3=\begin{pmatrix}1&0\\0&-1\end{pmatrix}.
\]
Bezpośrednie mnożenie daje
$\sigma_i^2=I_2$ i
$\sigma_i\sigma_j=-\sigma_j\sigma_i$ przy $i\ne j$;
na przykład $\sigma_1\sigma_2=i\sigma_3$.
Macierze $c(e_j)=i\sigma_j$ spełniają więc
\[
 c(e_i)c(e_j)+c(e_j)c(e_i)=-2\delta_{ij}I_2.
\]
Dla $n=2$ wystarczą $c(e_1),c(e_2)$.
Ich iloczyn to $-i\sigma_3$, a cztery macierze
$I_2,i\sigma_1,i\sigma_2,-i\sigma_3$
są liniowo niezależne i rozpinają
$M_2(\C)$. Z wymiaru $4$ obu stron dostajemy
$\operatorname{Cl}^{\C}_2\cong M_2(\C)$.
Przestrzeń spinorów $\Delta_2=\C^2$
jest zwykłą przestrzenią kolumn długości $2$.
\end{przyklad}

\begin{twierdzenie}[Klasyfikacja zespolona i wymiar spinorów]
\label{thm:cliff-kompleks}
Dla $m\geq0$ istnieją izomorfizmy algebr
\begin{align}
 \operatorname{Cl}^{\C}_{2m}
   &\cong M_{2^m}(\C), \label{eq:cliff-complex-even}\\
 \operatorname{Cl}^{\C}_{2m+1}
   &\cong M_{2^m}(\C)\oplus M_{2^m}(\C).
   \label{eq:cliff-complex-odd}
\end{align}
Wobec tego nieredukowalny moduł zespolony
ma wymiar $2^{\lfloor n/2\rfloor}$.
Dla parzystego $n$ jest jeden taki moduł
z dokładnością do izomorfizmu, a dla
nieparzystego $n$ są dwa.
\end{twierdzenie}
\begin{proof}
W wymiarze $0$ algebra ma postać $\C$.
W wymiarze $1$ składa się z $a+be_1$ przy
$e_1^2=-1$, więc podstawienia $e_1=i$
i $e_1=-i$ dają
$\operatorname{Cl}^{\C}_1\cong\C\oplus\C$.
Przykład~\ref{prz:pauli} daje
$\operatorname{Cl}^{\C}_2\cong M_2(\C)$.

Wykażmy krok indukcyjny zwiększający wymiar o $2$.
Iloczyn tensorowy algebr wyposażamy tu w mnożenie
$(a\otimes A)(b\otimes B)=ab\otimes AB$, rozszerzone
dwuliniowo; własność uniwersalna iloczynu tensorowego zapewnia
jego poprawne określenie. Łączność sprawdzamy na tensorach
prostych, a jedynką jest $1\otimes I_2$.
Niech $e_1,\ldots,e_n$ będą generatorami
$\operatorname{Cl}^{\C}_n$. W algebrze
$\operatorname{Cl}^{\C}_n\otimes M_2(\C)$
weźmy $n+2$ elementów:
\[
 F_j=e_j\otimes\sigma_3\ (j\leq n),\qquad
 F_{n+1}=1\otimes i\sigma_1,\qquad
 F_{n+2}=1\otimes i\sigma_2.
\]
Każdy $F_j^2=-1$. Dwa różne elementy
antykomutują: dla $j,\ell\leq n$ wynika
to z relacji starych $e_j$; między pierwszymi
$n$ a ostatnimi dwoma wynika z tego, że
$\sigma_3$ antykomutuje z $\sigma_1,\sigma_2$;
ostatnie dwa antykomutują jako macierze Pauliego.
Własność uniwersalna daje homomorfizm
\[
 \operatorname{Cl}^{\C}_{n+2}
    \longrightarrow
   \operatorname{Cl}^{\C}_n\otimes M_2(\C).
\]
Obrazy dwóch ostatnich generatorów dają całą
podalgebrę $1\otimes M_2(\C)$:
ich iloczyn jest $-i(1\otimes\sigma_3)$,
a wraz z jedynką tworzą bazę czterech macierzy.
Następnie z
$(e_j\otimes\sigma_3)(1\otimes\sigma_3)
=e_j\otimes I_2$ otrzymujemy wszystkie
generatory pierwszego czynnika. Homomorfizm
jest więc surjektywny. Obie algebry mają
wymiar $2^{n+2}$ nad $\C$
(Twierdzenie~\ref{thm:cliff-baza}),
więc jest izomorfizmem. Używamy przy tym
$M_d(\C)\otimes M_2(\C)\cong M_{2d}(\C)$:
jednostka macierzowa $E_{ab}\otimes E_{rs}$ odpowiada
jednostce z indeksem wiersza $(a,r)$ i kolumny $(b,s)$.
To bijekcja baz zachowująca mnożenie. Iloczyn tensorowy
rozdziela się też względem sumy prostej obu składników.
Iteracja od $n=0$
i $n=1$ daje~\eqref{eq:cliff-complex-even}
i~\eqref{eq:cliff-complex-odd}.

Pozostaje uzasadnić wymiary modułów. Algebra
$M_d(\C)$ działa na $\C^d$ nieredukowalnie:
jeśli podprzestrzeń zawiera $x\ne0$, to
dla dowolnego $y\in\C^d$ można znaleźć
macierz $A$ z $Ax=y$, więc zawiera wszystkie $y$.
Każdy nieredukowalny moduł $M_d(\C)$ jest
izomorficzny z tym standardowym: wybieramy
wektor $\psi$ z $E_{11}\psi\ne0$ (w razie potrzeby
zmieniamy indeks; jest to możliwe, gdyż
$\sum_jE_{jj}=I$). Wektory
$E_{i1}\psi$, $1\leq i\leq d$, są liniowo
niezależne: zastosowanie $E_{1j}$ do ich
kombinacji liniowej zerowej pozostawia
tylko $j$-ty współczynnik razy
$E_{11}\psi\ne0$. Ich rozpiętość zachowują
wszystkie macierze, bo
$E_{ab}(E_{i1}\psi)=\delta_{bi}E_{a1}\psi$.
Jest zatem niezerowym podmodułem, a
nieredukowalność wymusza, że stanowi cały moduł.
Dla sumy dwóch algebr macierzowych jej
dwa centralne idempotenty $(I,0)$, $(0,I)$
rozkładają każdy moduł na dwie niezmiennicze
części. W module nieredukowalnym tylko jedna
z nich jest niezerowa, co daje dokładnie
dwa wybory modułu wymiaru $d$.
\end{proof}

\begin{uwaga}[Dwa moduły w wymiarze nieparzystym]
Gdy $n=2m+1\geq3$, dwa nieredukowalne moduły
$\operatorname{Cl}^{\C}_n$ są różne, lecz
po ograniczeniu do $\operatorname{Spin}(n)$
dają izomorficzne reprezentacje.
Rzeczywiście, podalgebra parzysta jest
izomorficzna z
$\operatorname{Cl}^{\C}_{n-1}$: przypiszmy
jego generatorowi $e_j$ element $e_je_n$
dla $j<n$. Te elementy mają kwadrat $-1$
i wzajemnie antykomutują, a ponadto
generują wszystkie parzyste słowa
(każde $e_je_\ell$ z $j,\ell<n$ jest
iloczynem odpowiednich $e_je_n$ i $e_\ell e_n$
z dokładnością do znaku). Obie strony mają
wymiar $2^{n-1}$, więc to izomorfizm.
Ponieważ $n-1=2m$, algebra parzysta jest
$M_{2^m}(\C)$ i ma tylko jeden nieredukowalny
moduł. Trzeba jeszcze sprawdzić nieredukowalność ograniczeń:
każdy z dwóch modułów ma wymiar $d=2^m$, a w dowolnym
niezerowym module $M_d(\C)$ konstrukcja z jednostkami $E_{ij}$
z poprzedniego dowodu daje podmoduł wymiaru $d$.
Tutaj musi on być całą przestrzenią. Ponadto słowa parzyste
z wektorów ortonormalnych należą do Spin i rozpinają algebrę
parzystą, więc niezmienniczość względem Spin jest równoważna
niezmienniczości względem tej algebry.
W szczególności spinor dla obrotów
wymiaru $3$ ma po prostu dwie zespolone składowe.
\end{uwaga}

\begin{stwierdzenie}[Spinory chiralne w wymiarze parzystym]
\label{prop:spin-chiral}
Dla $n=2m\geq2$ ustalmy orientację $V$ i dodatnią bazę
ortonormalną $(e_1,\ldots,e_{2m})$. W module $\Delta_{2m}$
element
\[
 \Gamma:=i^m c(e_1e_2\cdots e_{2m})
\]
ma własność $\Gamma^2=I$, komutuje z działaniem
$\operatorname{Spin}(2m)$ i dzieli przestrzeń
na dwa podmoduły dla grupy Spin:
\[
 \Delta_{2m}=\Delta^+_{2m}\oplus\Delta^-_{2m},
 \qquad
 \Delta^\pm_{2m}=\ker(\Gamma\mp I),\qquad
 \dim_\C\Delta^\pm_{2m}=2^{m-1}.
\]
Nazywamy je \emph{półspinorami}.
\end{stwierdzenie}
\begin{proof}
Niech $\omega=e_1\cdots e_{2m}$. Aby uporządkować czynniki
w iloczynie $\omega^2$, wykonujemy $(2m)(2m-1)/2$ transpozycji;
następnie mnożymy $2m$ kwadratów $e_j^2=-1$.
Otrzymujemy
$\omega^2=(-1)^{m(2m-1)}(-1)^{2m}=(-1)^m$.
Ponieważ $i^{2m}=(-1)^m$, mamy $\Gamma^2=I$.
Zatem $P_\pm=\tfrac12(I\pm\Gamma)$ są projekcjami
o sumie $I$ i zerowym iloczynie; dają podany rozkład na przestrzenie własne.
Zmiana dodatniej bazy ortonormalnej nie zmienia $\omega$:
jej iloczyn jest antysymetryzacją najwyższego iloczynu zewnętrznego,
a ten zmienia się o wyznacznik, równy $1$.
Przy odwróceniu orientacji $\Gamma$ zmienia znak i oznaczenia
$\Delta^+$, $\Delta^-$ zamieniają się miejscami.
Przesunięcie dowolnego $e_j$ przez pozostałe
$2m-1$ generatory zmienia znak:
$\Gamma c(e_j)=-c(e_j)\Gamma$.
Zatem $\Gamma$ komutuje z iloczynami
parzystej liczby generatorów, w szczególności
ze $\operatorname{Spin}(2m)$; obie przestrzenie
własne są niezmiennicze. Operator
$c(e_1)$ jest odwracalny i przestawia
te przestrzenie własne, więc mają równe
wymiary. Suma wymiarów to $2^m$ z poprzedniego
twierdzenia, co daje $2^{m-1}$.
\end{proof}

\begin{przyklad}[Półspinory w płaszczyźnie]
Dla $n=2$ w modelu Pauliego
$\Gamma=i\,c(e_1)c(e_2)=\sigma_3$,
więc $\Delta^+_2=\C\binom10$ oraz
$\Delta^-_2=\C\binom01$.
Element $s_\theta$ z Przykładu~\ref{prz:spin2}
działa przez
\[
 c(s_\theta)=
 \begin{pmatrix}
  e^{-i\theta/2}&0\\0&e^{i\theta/2}
 \end{pmatrix}.
\]
Wektor po obrocie o $2\pi$ wraca do swego
położenia, a każdy z obu półspinorów
zostaje przemnożony przez $-1$.
\end{przyklad}

\begin{przyklad}[Spinor w $\R^3$ i sfera Blocha]
\label{prz:bloch}
W wymiarze $3$ działanie
$c(e_j)=i\sigma_j$ na $\Delta_3=\C^2$
spełnia relacje Clifforda. Na parzystych
generatorach z Przykładu~\ref{prz:spin3}
otrzymujemy macierze $-i\sigma_1,-i\sigma_2,
-i\sigma_3$. Dla
$s=a+bI+cJ+dK\in\operatorname{Spin}(3)$
macierz działania ma postać
\[
 U_s=aI_2-i(b\sigma_1+c\sigma_2+d\sigma_3).
\]
Z relacji Pauliego dostajemy
$U_s^*U_s=(a^2+b^2+c^2+d^2)I_2=I_2$
oraz $\det U_s=1$. Każda macierz z
$\mathrm{SU}(2)$ ma taką postać: z warunków
unitarności i wyznacznika $1$ wynika
$U=\begin{psmallmatrix}\alpha&\beta\\
-\overline\beta&\overline\alpha\end{psmallmatrix}$,
gdzie $|\alpha|^2+|\beta|^2=1$; cztery
rzeczywiste części $\alpha,\beta$ odpowiadają
jedynej czwórce $(a,b,c,d)$.
Zatem $\operatorname{Spin}(3)\cong\mathrm{SU}(2)$.

Niech teraz $\psi=(z_1,z_2)^T$ ma normę $1$.
Liczby rzeczywiste
\begin{equation}\label{eq:spin-bloch}
 b_j(\psi):=\psi^*\sigma_j\psi
 \quad\text{dają}\quad
 b(\psi)=
 \bigl(2\operatorname{Re}(\overline z_1z_2),
       2\operatorname{Im}(\overline z_1z_2),
       |z_1|^2-|z_2|^2\bigr).
\end{equation}
Rachunek daje
$\|b(\psi)\|^2=
4|z_1|^2|z_2|^2+(|z_1|^2-|z_2|^2)^2=1$.
Zmiana fazy $\psi\mapsto e^{i\varphi}\psi$
nie zmienia $b(\psi)$. Odwrotnie, liczby
$b_j$ wyznaczają macierz
$\psi\psi^*=\frac12(I_2+\sum_jb_j\sigma_j)$,
a ta macierz wyznacza prostą $\C\psi$.
Każdy punkt sfery otrzymujemy: dla
$b=(x,y,z)$ z $z>-1$ bierzemy
$z_1=\sqrt{(1+z)/2}$,
$z_2=(x+iy)/\sqrt{2(1+z)}$;
dla bieguna południowego bierzemy $(0,1)$.
Otrzymujemy więc konkretną identyfikację
\[
 \mathbb{CP}^1=\operatorname{Gr}_1(\C^2)\cong S^2.
\]
Jest to także gładka identyfikacja: na mapie $[1:w]$
wzór ma postać
$b=(2\operatorname{Re}w,2\operatorname{Im}w,1-|w|^2)/(1+|w|^2)$,
a odwrotność poza biegunem południowym to $w=(x+iy)/(1+z)$;
druga mapa analogicznie obsługuje biegun południowy.
Macierze $U_s$ obracają odpowiedni punkt sfery.
Istotnie, z $U_sc(v)U_s^{-1}=c(\rho(s)v)$ i $c(v)=i\sum_jv_j\sigma_j$
otrzymujemy
\[
 (U_s\psi)(U_s\psi)^*
 =U_s\psi\psi^*U_s^{-1}
 =\tfrac12\bigl(I_2+\textstyle\sum_j(\rho(s)b(\psi))_j\sigma_j\bigr).
\]
Jednoznaczność współczynników przy $\sigma_j$ daje
$b(U_s\psi)=\rho(s)b(\psi)$.
Znak spinora $\psi\mapsto-\psi$ jest
niewidoczny po przejściu do punktu sfery.
\end{przyklad}

\begin{figure}[!htbp]
\centering
\begin{tikzpicture}[>=Stealth,line cap=round,scale=1.0]
 \draw[manifoldblue!75!black,thick,rounded corners]
 (-5.6,-.72) rectangle (-2.62,.94);
 \node[font=\small] at (-4.11,.49)
   {$\psi=\binom{z_1}{z_2}$};
 \node[font=\small] at (-4.11,-.17)
   {$|z_1|^2+|z_2|^2=1$};
 \draw[->,manifoldblue,thick] (-3.65,1.17)
  arc[start angle=0,end angle=300,x radius=.44,y radius=.18];
 \node[font=\scriptsize,manifoldblue] at (-4.42,1.47)
  {fazy $e^{i\varphi}\psi$};
 \draw[->,thick,manifoldgold!85!black] (-2.4,.12)--(-.85,.12)
  node[midway,above,font=\small] {$\psi\mapsto\C\psi$};
 \node[draw=manifoldgreen!80!black,rounded corners,
  fill=manifoldgreen!8,font=\small,inner sep=6pt] at (.05,.12)
  {$[\psi]\in\mathbb{CP}^1$};
 \draw[->,thick,manifoldgold!85!black] (.97,.12)--(2.1,.12)
  node[midway,above,font=\small] {$b$};
 \shade[inner color=white,outer color=manifoldblue!43] (3.23,.12) circle (.98);
 \draw[manifoldblue!75!black,thick] (3.23,.12) circle (.98);
 \draw[manifoldgreen!55,thin] (2.25,.12)
   arc[start angle=180,end angle=360,x radius=.98,y radius=.23];
 \draw[manifoldgreen!55,thin,dashed] (4.21,.12)
   arc[start angle=0,end angle=180,x radius=.98,y radius=.23];
 \fill[manifoldred] (3.65,.66) circle (2.3pt);
 \node[manifoldred,font=\small] at (4.56,.85) {$b(\psi)\in S^2$};
\end{tikzpicture}
\caption{Spinor jednostkowy $\psi\in\C^2$ ma fazę,
która nie zmienia prostej $[\psi]\in\mathbb{CP}^1$.
Wzór~\eqref{eq:spin-bloch} przypisuje tej prostej
punkt sfery. To łączy spinory z grassmannianem prostych.}
\label{fig:spin-bloch}
\end{figure}

\section{Struktury spinowe i wiązki spinorów}
\label{sec:struktura-spinowa}

Dotąd wektory, obroty i spinory należały do jednej
ustalonej przestrzeni $\R^n$. Na rozmaitości
przestrzeń $T_pM$ zmienia się z $p$.
Gładka rodzina dodatnio określonych iloczynów
skalarnych $g_p$ na $T_pM$ to \emph{metryka
riemannowska}. Pozwala wybierać bazy ortonormalne
w każdym punkcie. Jeśli $M$ jest zorientowana,
możemy wybierać bazy o dodatniej orientacji.

\emph{Wiązka dodatnio zorientowanych ramek
ortonormalnych} $F^+(M)\to M$ ma w punkcie $p$
za włókno wszystkie uporządkowane bazy
$(u_1,\ldots,u_n)$ przestrzeni $T_pM$,
które są ortonormalne i dodatnie.
Grupa $\mathrm{SO}(n)$ działa z prawej strony
przez zmianę takiej bazy. Nie ma wyróżnionej
bazy w każdym włóknie; lokalnie wiązka wygląda
jednak jak $U\times\mathrm{SO}(n)$.
Wiązkę o takiej lokalnej postaci, z wolnym
działaniem grupy przechodzącym tranzytywnie
na każdym włóknie, nazywamy \emph{wiązką główną}.

\begin{definicja}[Struktura spinowa i wiązka spinorów]
\index{struktura spinowa i wiazka spinorow@Struktura spinowa i wiązka spinorów}
\label{def:struktura-spinowa}
Dla zorientowanej rozmaitości riemannowskiej
$M^n$, $n\geq2$, \emph{struktura spinowa}
składa się z wiązki głównej
$P_{\mathrm{Spin}}\to M$ o grupie
$\operatorname{Spin}(n)$ oraz odwzorowania
gładkiego $\lambda:P_{\mathrm{Spin}}\to F^+(M)$
nad identycznością bazy, takiego że
\[
 \lambda(pg)=\lambda(p)\rho(g),
 \qquad g\in\operatorname{Spin}(n),
\]
a $\lambda$ ograniczona do każdego włókna
jest dwukrotnym nakryciem.
\emph{Wiązka spinorów} to wiązka zespolona
\[
 S:=P_{\mathrm{Spin}}\times_{\operatorname{Spin}(n)}\Delta_n.
\]
Jej elementami są klasy par $(p,\psi)$,
$p\in P_{\mathrm{Spin}}$, $\psi\in\Delta_n$,
z relacją $(pg,\psi)\sim(p,c(g)\psi)$.
\emph{Polem spinorowym} nazywamy gładki przekrój $S$.
Tę samą definicję stosujemy do dowolnej zorientowanej wiązki
euklidesowej $E$ rzędu $n$, zastępując $F^+(M)$ wiązką jej
dodatnich ramek ortonormalnych $F^+(E)$.
\end{definicja}

Po wyborze lokalnej ramki spinowej $s_i:U_i\to P_{\mathrm{Spin}}$
wiązka $S$ wygląda jak $U_i\times\Delta_n$. Ustalmy kierunek
zmian: jeśli $s_j=s_i h_{ij}$, to
$[s_j,\psi_j]=[s_i,c(h_{ij})\psi_j]$, zatem współrzędne
tego samego spinora spełniają $\psi_i=c(h_{ij})\psi_j$.
Przy przejściu od kolumny $i$ do kolumny $j$ używamy więc
macierzy odwrotnej. Z równości dla lokalnych ramek wynika też
$h_{ij}h_{jk}=h_{ik}$; reprezentacja $c$ przenosi tę równość
na macierze przejścia, co uzasadnia strukturę wiązki wektorowej $S$.
Zauważmy też, że $\lambda$ jest nakryciem całych przestrzeni:
w lokalnej trywializacji zadanej przez $s_i$ i ramkę $\lambda\circ s_i$
ma postać $(x,g)\mapsto(x,\rho(g))$.
Jeśli zmiana ramki wektorowej wykonuje pełny
obrót o $2\pi$, jej podniesienie w Spin
może dojść do $-1$. Wtedy wektory wracają
do wartości początkowych, lecz spinor dostaje
znak minus. To lokalna treść różnicy między
wiązką wektorów i wiązką spinorów.

\begin{przyklad}[Płaszczyzna i sfera są spinowe]
Na $\R^n$ standardowe ramki dają
$F^+(\R^n)\cong\R^n\times\mathrm{SO}(n)$.
Bierzemy
$P_{\mathrm{Spin}}=\R^n\times\operatorname{Spin}(n)$
i $\lambda(x,s)=(x,\rho(s))$; spinory
tworzą trywialną wiązkę
$S\cong\R^n\times\Delta_n$.

Na sferze $S^2\subset\R^3$ każda dodatnia
ramka styczna $(u,v)$ w punkcie $p$ wyznacza
macierz o kolumnach $(u,v,p)\in\mathrm{SO}(3)$.
To utożsamia $F^+(S^2)$ z $\mathrm{SO}(3)$,
a rzut do bazy wysyła macierz na jej trzecią
kolumnę $p$. Dwukrotne nakrycie
$\mathrm{SU}(2)\cong\operatorname{Spin}(3)
\to\mathrm{SO}(3)$ z
Przykładu~\ref{prz:spin3} daje odwzorowanie $\lambda$.
Rzut z $\operatorname{Spin}(3)$ do sfery ma postać
$q(s)=\rho(s)e_3$. Używamy jednej, ustalonej podgrupy
$H=\rho^{-1}(\operatorname{Stab}(e_3))=\operatorname{Spin}(2)$,
złożonej z $\cos t+\sin t\,e_1e_2$. Działa ona z prawej strony
na $\operatorname{Spin}(3)$ i zachowuje $q(s)$.
Jeśli $q(s)=q(s')$, to $s^{-1}s'\in H$, więc działanie jest
wolne i przechodnie na każdym włóknie. Lokalne gładkie ramki
styczne $(u,v)$ dają lokalne sekcje w $\mathrm{SO}(3)$;
po zmniejszeniu otoczenia podnosimy je przez lokalną odwrotność
nakrycia $\rho$. To zapewnia lokalną trywialność wiązki głównej.
Równość $\lambda(sh)=\lambda(s)\rho(h)$ kończy sprawdzenie
definicji struktury spinowej na sferze.
\end{przyklad}

\begin{uwaga}[Podnoszenie na zbiorze ściągalnym]
\label{rem:spin-podnoszenie}
Użyjemy elementarnego faktu o nakryciach: gładkie odwzorowanie
ze ściągalnego obszaru $U$ do podstawy gładkiego nakrycia ma
gładkie podniesienie po wybraniu obrazu jednego punktu.
Istotnie, drogę od tego punktu podnosimy kolejno przez lokalne
odwrotności nakrycia, dzieląc jej przedział parametrów na skończenie
wiele odcinków. Jednoznaczność na przecięciach daje jednoznaczność
podniesionej drogi. Homotopię dróg o ustalonych końcach podnosimy
analogicznie po podziale kwadratu parametrów na małe prostokąty,
których obrazy leżą w dziedzinach lokalnych odwrotności.
Końcowy punkt podniesienia nie zmienia się w tej homotopii:
poruszałby się ciągle w dyskretnym włóknie. W obszarze ściągalnym
każde dwie drogi o tych samych końcach są tak homotopijne,
więc ich podniesienia wyznaczają tę samą wartość w punkcie końcowym.
Otrzymana funkcja jest lokalnie złożeniem danej funkcji z jedną
lokalną odwrotnością nakrycia, zatem jest gładka.
Ten sam argument działa na dysku z brzegiem.
\end{uwaga}

\begin{przyklad}[Wiązka, której nie da się podnieść]
\label{prz:nie-spin-linia}
Podłoże $S^2\cong\mathbb{CP}^1$ jest spinowe,
ale nie każda \emph{wiązka} nad nim dopuszcza
podniesienie struktury do Spin.
Weźmy zespoloną wiązkę tautologiczną
$\gamma^1_\C$ nad
$\operatorname{Gr}_1(\C^2)=\mathbb{CP}^1$
i potraktujmy jej włókna jako rzeczywiste
zorientowane płaszczyzny wymiaru $2$.
Na mapie $[1:z]$ mamy ramkę zespoloną
$s_0(z)=(1,z)$, a na mapie $[w:1]$,
$w=1/z$, ramkę $s_\infty(w)=(w,1)$.
Na przecięciu
$s_\infty(1/z)=z^{-1}s_0(z)$.
Po unormowaniu ramek, przy $|z|=1$, $z=e^{i\theta}$, unitarny
czynnik przejścia jest więc $e^{-i\theta}$:
jako rzeczywista izometria obraca płaszczyznę
o kąt $-\theta$. Po przejściu
$0\leq\theta\leq2\pi$ daje pętlę
o nieparzystej liczbie pełnych obrotów.
Ta pętla w $\mathrm{SO}(2)$ nie ma
\emph{zamkniętego} podniesienia do $\operatorname{Spin}(2)$:
połowa kąta biegnie
od $0$ do $-\pi$ i kończy w $-1$, nie w $1$.
Gdyby istniała struktura spinowa, ramkę ortonormalną nad każdą
z dwóch półkul można by podnieść do ramki spinowej.
Rzeczywiście, przeciwobraz takiej ramki przez $\lambda$ jest
dwukrotnym nakryciem dysku; Uwagi~\ref{rem:spin-podnoszenie}
używamy do podniesienia identyczności dysku. Przejście między
otrzymanymi ramkami spinowymi na równiku byłoby właśnie
zamkniętym podniesieniem opisanej pętli, co jest sprzecznością.
Zmiana ramki na którejś z dwóch półkul
zmienia pętlę tylko o taką, która przedłuża
się na dysk i ma zerowy całkowity obrót.
Dlatego ten brak podniesienia nie zależy
od wybranych ramek. Wiązka rzeczywista otrzymana przez zapomnienie
struktury zespolonej $\gamma^1_\C$ nie ma struktury
spinowej, mimo że sama sfera $S^2$ ją ma.
\end{przyklad}

\begin{stwierdzenie}[Znaki na potrójnych przecięciach]
\label{prop:w2-znaki}
Niech $E$ będzie zorientowaną wiązką euklidesową
rzędu $n\geq2$ nad $M$. Weźmy dostatecznie drobne
pokrycie trywializujące $(U_i)$, w którym
niepuste skończone przecięcia są ściągalne. Takie pokrycie
nazywamy \emph{dobrym}; w tym stwierdzeniu zakładamy jego wybór.
Niech $g_{ij}:U_i\cap U_j\to\mathrm{SO}(n)$
przepisują współrzędne z trywializacji $j$ do trywializacji $i$,
więc $g_{ij}g_{jk}=g_{ik}$.
Na przecięciach wybierzmy podniesienia
$h_{ij}:U_i\cap U_j\to\operatorname{Spin}(n)$
z $\rho(h_{ij})=g_{ij}$ i $h_{ji}=h_{ij}^{-1}$.
Wtedy na potrójnym przecięciu
\begin{equation}\label{eq:w2-znaki}
 \varepsilon_{ijk}:=h_{ij}h_{jk}h_{ki}
          \in\{1,-1\}.
\end{equation}
Wiązka $E$ ma strukturę spinową dokładnie wtedy,
gdy wybory $h_{ij}$ można zmodyfikować znakami
tak, aby wszystkie $\varepsilon_{ijk}$ wynosiły $1$.
\end{stwierdzenie}
\begin{proof}
Z warunku kocyklu dla $g_{ij}$ mamy
$\rho(\varepsilon_{ijk})
=g_{ij}g_{jk}g_{ki}=I$.
Twierdzenie~\ref{thm:spin-nakrycie} mówi, że
jądro $\rho$ to $\{\pm1\}$; ponieważ znak jest
dyskretny, $\varepsilon_{ijk}$ jest lokalnie stały.
Nie ma problemu z samym podniesieniem
$g_{ij}$ na ściągalnym przecięciu:
zapewnia je Uwaga~\ref{rem:spin-podnoszenie}.
Przyjmujemy też $h_{ii}=1$.

Jeżeli na każdym potrójnym przecięciu
$\varepsilon_{ijk}=1$, to
$h_{ij}h_{jk}=h_{ik}$.
Są to funkcje przejścia z warunkiem kocyklu.
Sklejamy $U_i\times\operatorname{Spin}(n)$ relacją
$(x,a)_j\sim(x,h_{ij}(x)a)_i$. Warunek kocyklu zapewnia
przechodniość tej relacji, a gładkie przejścia zapewniają lokalne
trywializacje, jak w Sekcji~\ref{sec:wiazki-wektorowe}.
Działanie z prawej strony jest dobrze określone, wolne
i przechodnie we włóknach, więc otrzymujemy wiązkę główną.
Jeśli $f_i$ oznacza wybraną lokalną ramkę $E$, to
$\lambda([x,a]_i)=f_i(x)\rho(a)$ jest niezależne od reprezentanta,
ponieważ $f_j=f_i g_{ij}$ i $\rho(h_{ij})=g_{ij}$.
Lokalnie jest to $\id\times\rho$, a więc spełnia definicję
struktury spinowej.
Odwrotnie, w istniejącej strukturze spinowej podnosimy każdą
ustaloną ramkę $f_i$ do ramki spinowej nad ściągalnym $U_i$.
Stosujemy Uwagę~\ref{rem:spin-podnoszenie} do dwukrotnego
nakrycia będącego przeciwobrazem $f_i$ przez $\lambda$.
Przejścia tych ramek podnoszą właśnie dane $g_{ij}$ i spełniają
warunek kocyklu, więc wszystkie znaki są $1$.
Dowolne dwa odwzorowania $h_{ij}$ nad
tym samym $g_{ij}$ różnią się na spójnym
przecięciu stałym znakiem
$\delta_{ij}\in\{\pm1\}$.
Zmiana podniesień zastępuje
$\varepsilon_{ijk}$ przez
$\delta_{ij}\delta_{jk}\delta_{ki}
\varepsilon_{ijk}$, co dowodzi
sformułowania o modyfikacji znakami.
\end{proof}

\begin{uwaga}[Druga klasa Stiefela--Whitneya]
Zdarza się, że pojedynczych znaków
$\varepsilon_{ijk}$ nie da się wszystkich
usunąć. Ich układ spełnia warunek zgodności
na poczwórnych przecięciach: dwa sposoby
pomnożenia $h_{ij}h_{jk}h_{k\ell}$
dają
$\varepsilon_{ijk}\varepsilon_{ik\ell}
=\varepsilon_{jk\ell}\varepsilon_{ij\ell}$.
Taką rodzinę znaków nazywamy
\emph{dwukocyklowym} zapisem przeszkody.
Zmiany przez $\delta_{ij}$ z dowodu
uznajemy za równoważne. Nie zmienia tego wybór lokalnych ramek:
jeśli $f'_i=f_i a_i$, wybieramy podniesienia $b_i$ funkcji $a_i$
nad ściągalnymi $U_i$. Wtedy $g'_{ij}=a_i^{-1}g_{ij}a_j$
mają podniesienia $h'_{ij}=b_i^{-1}h_{ij}b_j$.
Ich iloczyn na potrójnym przecięciu to
$b_i^{-1}\varepsilon_{ijk}b_i=\varepsilon_{ijk}$,
bo znaki należą do centrum grupy Spin.
Klasa tego układu w
$H^2(M;\Z/2)$ nazywa się
\emph{drugą klasą Stiefela--Whitneya}
$w_2(E)$. Na dobrym pokryciu
$w_2(E)=0$ oznacza dokładnie, że znaki
można usunąć, czyli że struktura spinowa
istnieje. Identyfikacja klas z różnych
pokryć to twierdzenie kohomologii Čech;
używamy tu tylko elementarnego testu
podnoszenia funkcji przejścia dowiedzionego
wyżej. Dla wiązki w
Przykładzie~\ref{prz:nie-spin-linia}
przeszkoda jest niezerowa.
\end{uwaga}

\par\smallskip
{\footnotesize\noindent\emph{Dalsza lektura:}
Rachunki z konwencją $v^2=-\|v\|^2$
można porównać z notatkami Briana Conrada,
\emph{Clifford algebras and spin groups},
\url{https://math.stanford.edu/~conrad/210CPage/handouts/clifford.pdf}.
Budowę struktur spinowych i opis przeszkody
$w_2$ rozwijają notatki Konstantina Wernliego,
\emph{Lecture Notes on Spin Geometry},
\url{https://arxiv.org/pdf/1911.09766}.
Obie lektury używają miejscami innej kolejności tematów.
Powyższy test znaków został dowiedziony bez użycia kohomologii;
jego interpretacja jako klasy w $H^2$ korzysta z teorii
kohomologii Čech, której tutaj nie dowodziliśmy.\par}
\clearpage
